\documentclass{article}

% if you need to pass options to natbib, use, e.g.:
%     \PassOptionsToPackage{numbers, compress}{natbib}
% before loading neurips_2026

% The authors should use one of these tracks.
% Before accepting by the NeurIPS conference, select one of the options below.
% "default" for submission
% \usepackage{neurips_2026_vericode}
% Camera-ready: the .sty's "final" option.
\usepackage[final]{neurips_2026_vericode}

% After being accepted, the authors should add "final" behind the track to compile a camera-ready version.
% \usepackage[main, final]{neurips_2026_vericode}

% "preprint" option is used for arXiv or other preprint submissions
 % \usepackage[preprint]{neurips_2026_vericode}

% to avoid loading the natbib package, add option nonatbib:
%    \usepackage[nonatbib]{neurips_2026_vericode}

\usepackage[utf8]{inputenc} % allow utf-8 input
\usepackage[T1]{fontenc}    % use 8-bit T1 fonts
\usepackage{hyperref}       % hyperlinks
\usepackage{url}            % simple URL typesetting
\usepackage{booktabs}       % professional-quality tables
\usepackage{amsfonts}       % blackboard math symbols
\usepackage{nicefrac}       % compact symbols for 1/2, etc.
\usepackage{microtype}      % microtypography
\usepackage{xcolor}         % colors


% Figures live in macros, never as literals in prose (paper-vericode/CLAUDE.md).
% numbers.tex - every reported figure in this paper, as a macro.
%
% Rule (paper-vericode/CLAUDE.md): no literal number appears in prose.
% Every macro below is transcribed from a committed CSV under ../results/.
% Each comment names the source file, the row selector, and the column.
% Macros ending in "Pct" render a percentage: either 100x a named fraction
% column, or a figure already recorded as a percentage and transcribed as one.
% Each macro's own comment says which. Every other transformation in this file is arithmetic on
% committed values and prints its own formula in the macro's comment: the
% derived block near the end (quotients) and \stepOverlapCensus (a column sum).

% ---------------------------------------------------------------------------
% Corpus sizes
% ---------------------------------------------------------------------------
\newcommand{\humanProofs}{70,086}      % results/coverage_common_support.csv, scope=full, human_proofs
\newcommand{\machineProofs}{29,750}    % results/coverage_common_support.csv, scope=full, machine_proofs
\newcommand{\humanProofsSupport}{70,013} % results/coverage_common_support.csv, scope=common_support, human_proofs
\newcommand{\matchedProofs}{22,785}    % results/coverage_common_support.csv, scope=matched_1to1, machine_proofs

% ---------------------------------------------------------------------------
% Coverage - full corpora, corrected premise sets (H2a headline)
% The uncorrected `raw` policy row (118,514 / 0.0089 / 99.1%) is retracted and
% is deliberately not given a macro.
% ---------------------------------------------------------------------------
\newcommand{\humanVocab}{78,642}       % results/coverage_by_policy.csv, policy=corrected, human_vocab
\newcommand{\machineVocab}{1,127}    % results/coverage_by_policy.csv, policy=corrected, machine_vocab
\newcommand{\vocabRatio}{69.8}         % results/coverage_area_restricted.csv, design=full corpora, ratio (69.77994676131323)
\newcommand{\jaccardFull}{0.0129}      % results/coverage_by_policy.csv, policy=corrected, jaccard (0.012875372992190971)
\newcommand{\intersectionFull}{1,014} % results/coverage_by_policy.csv, policy=corrected, intersection
\newcommand{\humanOnlyFullPct}{98.7}   % results/coverage_by_policy.csv, policy=corrected, human_only_fraction (0.9871061264972916), x100
\newcommand{\machineOnlyFull}{113}     % results/coverage_by_policy.csv, policy=corrected, machine_only

% ---------------------------------------------------------------------------
% Coverage - common support and 1:1 matched designs
% ---------------------------------------------------------------------------
\newcommand{\humanVocabSupport}{78,591} % results/coverage_common_support.csv, scope=common_support, human_vocab
\newcommand{\machineVocabSupport}{1,123} % results/coverage_common_support.csv, scope=common_support, machine_vocab
\newcommand{\jaccardSupport}{0.0128}   % results/coverage_common_support.csv, scope=common_support, jaccard (0.012820024140778858)
\newcommand{\intersectionSupport}{1,009} % results/coverage_common_support.csv, scope=common_support, intersection
\newcommand{\humanOnlySupportPct}{98.7} % results/coverage_common_support.csv, scope=common_support, human_only_fraction (0.9871613798017584), x100
\newcommand{\humanVocabMatched}{21,162}  % results/coverage_common_support.csv, scope=matched_1to1, human_vocab
\newcommand{\machineVocabMatched}{1,123} % results/coverage_common_support.csv, scope=matched_1to1, machine_vocab
\newcommand{\jaccardMatched}{0.0360}   % results/coverage_common_support.csv, scope=matched_1to1, jaccard (0.03598159081400214)
\newcommand{\intersectionMatched}{774} % results/coverage_common_support.csv, scope=matched_1to1, intersection
\newcommand{\humanOnlyMatchedPct}{96.3} % results/coverage_common_support.csv, scope=matched_1to1, human_only_fraction (0.9634250070881769), x100

% ---------------------------------------------------------------------------
% Coverage - corpus-size and region controls
% ---------------------------------------------------------------------------
\newcommand{\humanVocabRarefied}{49,601} % results/vocabulary_rarefaction.csv, corpus=human fraction=0.4244784978455041 (proofs=29,750), vocab_mean
\newcommand{\rarefiedFivePctProofs}{3,504} % results/vocabulary_rarefaction.csv, corpus=human fraction=0.05, proofs
\newcommand{\humanVocabRarefiedFivePct}{11,493} % results/vocabulary_rarefaction.csv, corpus=human fraction=0.05, vocab_mean (11493.4)
\newcommand{\humanVocabAreaRestricted}{32,091} % results/coverage_area_restricted.csv, design=restricted - support (PRIMARY), human_vocab
\newcommand{\areaRestrictedRatio}{28.5} % results/coverage_area_restricted.csv, design=restricted - support (PRIMARY), ratio (28.474711623779946)
\newcommand{\areaRestrictedJaccard}{0.0315} % results/coverage_area_restricted.csv, design=restricted - support (PRIMARY), jaccard (0.03148677182958639)
\newcommand{\intersectionAreaRestricted}{1,014} % results/coverage_area_restricted.csv, design=restricted - support (PRIMARY), intersection; equal to \intersectionFull because the restriction drops no machine premise
\newcommand{\areaRestrictedHumanOnlyPct}{96.8} % results/coverage_area_restricted.csv, design=restricted - support (PRIMARY), human_only_fraction (0.9684023558006918), x100
\newcommand{\areaRegionsMachine}{42}    % results/coverage_area_restricted.csv, design=restricted - support (PRIMARY), n_regions
\newcommand{\areaRegionsHuman}{1,595}  % results/coverage_area_restricted.csv, design=full corpora, n_regions
\newcommand{\areaHumanSurvivalPct}{40.8} % results/coverage_area_restricted.csv, design=restricted - support (PRIMARY), human_survival (0.40806439307240405), x100
% Registered sensitivity S1. Reported because its machine leakage is the
% reason the support rule, not the top-12 decomposition, is primary.
\newcommand{\areaTopTwelveRatio}{19.2} % results/coverage_area_restricted.csv, design=S1 restricted - top-12 only, ratio (19.248071979434446)
\newcommand{\areaTopTwelveMachineLeakPct}{31.0} % results/coverage_area_restricted.csv, design=S1 restricted - top-12 only, machine_leak (0.30967169476486245), x100

% ---------------------------------------------------------------------------
% Extraction instrument - retention decomposition (mathematical layer, loose matcher)
% ---------------------------------------------------------------------------
\newcommand{\ceilingWholeFile}{0.7915} % results/extraction_audit.csv, layer=gate matcher=efficiency_paired, total
\newcommand{\reachableBound}{0.6948}   % results/extraction_audit.csv, layer=gate matcher=efficiency_reachable, total
\newcommand{\retention}{0.6124}        % results/extraction_audit.csv, layer=mathematical matcher=loose, rate
\newcommand{\efficiencyReachable}{0.8815} % results/extraction_audit.csv, layer=gate matcher=efficiency_reachable, rate

% ---------------------------------------------------------------------------
% Instrument reference line (Section 3.4).
%
% WHICH FILE CARRIES THE REFERENCE LINE. The frozen line is the STEP-AXIS
% JOINT value and it lives in results/length_controlled_reference.csv
% (src/config.py: H2_INSTRUMENT_REFERENCE = 0.099, CI (0.068, 0.126),
% band (0.076, 0.126), H2_INSTRUMENT_REFERENCE_AXIS = "steps").
% results/instrument_reference.csv is NOT that file: it carries the earlier,
% UNCONTROLLED class-only gap on a different sample (n=1,679 class-targeted,
% against n=7,531 for the standardization). It is deliberately not macro'd
% here -- a labelled macro for a superseded value is still a value a later
% session can quote by reaching for the obvious-sounding name.
% ---------------------------------------------------------------------------

% Recall term -- the reference line proper. Joint standardization on
% (tactic class x length bin), step axis.
\newcommand{\recallTermClassOnly}{0.0748} % results/length_controlled_reference.csv, axis=steps, gap_class_only (0.07479504334853565)
\newcommand{\recallTermJoint}{0.0986}     % results/length_controlled_reference.csv, axis=steps, gap_joint (0.09855205523992183)
\newcommand{\recallTermCIlow}{0.0678}     % results/length_controlled_reference.csv, axis=steps, ci_low (0.0678314769103629)
\newcommand{\recallTermCIhigh}{0.1256}    % results/length_controlled_reference.csv, axis=steps, ci_high (0.12564349773213496)
\newcommand{\recallTermBandLow}{0.0762}   % results/length_controlled_reference.csv, axis=steps, band_low (0.07619722282459518); 24 precedence orderings
\newcommand{\recallTermBandHigh}{0.1257}  % results/length_controlled_reference.csv, axis=steps, band_high (0.1256745904744775); 24 precedence orderings

% Precision term, both axes. Negative: the machine side carries the LOWER
% spurious rate, so the term displaces the comparison the same way the recall
% term does rather than cancelling it. Wrapped in \ensuremath so the minus
% sign sets as a minus and not a hyphen.
\newcommand{\spuriousStepAxis}{\ensuremath{-0.0142}} % results/spurious_length_controlled.csv, axis=steps, gap_joint (-0.014167715009446358)
\newcommand{\spuriousCharAxis}{\ensuremath{-0.1917}} % results/spurious_length_controlled.csv, axis=chars, gap_joint (-0.1917329526614106)

% The diagnostic that disqualifies the step-axis precision term: the share of
% the machine mix sitting in cells with no usable human rate. Dominated by the
% arith x 0-2 cell (\degenerateMachineSharePct{}% of the corpus) but not
% identical to it -- two thin decide cells carry the rest.
\newcommand{\spuriousCollapsedMachineMixPct}{24.1} % results/spurious_length_controlled.csv, axis=steps, collapsed_machine_mix (0.2408739495798319), x100

% ---------------------------------------------------------------------------
% Length overlap between the corpora, by axis
%
% TWO step-axis values exist and they are not interchangeable. The census
% figure is the one the paper quotes; the sample figure is what the
% standardization was actually cut on. Section 3.5 carries a footnote naming
% both, so both are macro'd rather than one being silently preferred.
% The char figure is on the sample basis, hence \charOverlapSample: the name
% says which basis, as it does for the two step-axis macros above.
% ---------------------------------------------------------------------------
\newcommand{\stepOverlapCensus}{0.8563} % results/length_overlap.csv, axis=steps, sum of overlap_mass over all six bins (0.26077310924369745 + 0.2699831932773109 + 0.18574322974631166 + 0.0994207116970579 + 0.034743029991724454 + 0.00561344537815126 = 0.8562767193342535); full census, both corpora
\newcommand{\stepOverlapSample}{0.8599} % results/length_controlled_reference.csv, axis=steps, overlap_coefficient (0.8599090577074936); n=7,531 standardization sample, NOT the census
\newcommand{\charOverlapSample}{0.2151} % results/length_controlled_reference.csv, axis=chars, overlap_coefficient (0.21511587708450383)

% ---------------------------------------------------------------------------
% Automation stratum shares (string matcher, censused mixes)
% ---------------------------------------------------------------------------
\newcommand{\humanArithMixPct}{0.7}    % results/tactic_class_mix.csv, corpus=human_mathlib instrument=matcher, arith (0.006777387780726536), x100
\newcommand{\machineArithMixPct}{76.4} % results/tactic_class_mix.csv, corpus=machine_goedel instrument=matcher, arith (0.7638655462184873), x100
\newcommand{\humanNoneMixPct}{53.4}    % results/tactic_class_mix.csv, corpus=human_mathlib instrument=matcher, none (0.5341437662300602), x100
\newcommand{\machineNoneMixPct}{4.0}   % results/tactic_class_mix.csv, corpus=machine_goedel instrument=matcher, none (0.04016806722689076), x100

% ---------------------------------------------------------------------------
% H2c - the degenerate cell (arith x 0-2 tactic steps)
% ---------------------------------------------------------------------------
\newcommand{\degenerateMachineSharePct}{23.1} % results/common_support_cells.csv, cell=arith x 0-2, machine_share (0.230890756302521), x100
\newcommand{\degenerateHumanSharePct}{0.02}   % results/common_support_cells.csv, cell=arith x 0-2, human_share (0.00019975458722141368), x100
\newcommand{\degenerateHumanCount}{14}        % results/common_support_cells.csv, cell=arith x 0-2, human_n
\newcommand{\degenerateMachineCount}{6,869} % results/common_support_cells.csv, cell=arith x 0-2, machine_n

% ---------------------------------------------------------------------------
% The published graph (Section 2). These are the RELEASE's own figures, not
% our measurements: data/mathlib_graph/v2/summary.json, fetched at the pinned
% revision by scripts/fetch_mathlib_graph.sh and published with
% arXiv:2604.24797. Not under results/ because we did not compute them; they
% are reproducible by fetching the release.
% ---------------------------------------------------------------------------
\newcommand{\graphDeclarations}{308,129}  % summary.json, declaration_graph.node_count
\newcommand{\graphEdges}{8,436,366}       % summary.json, declaration_graph.edge_count_raw
\newcommand{\explicitEdgeCount}{2,178,534} % summary.json, declaration_graph.explicit_edge_count
\newcommand{\explicitEdgePct}{25.8}       % summary.json, declaration_graph.explicit_edge_ratio (0.2582), x100
\newcommand{\synthesizedEdgePct}{74.2}    % summary.json, declaration_graph.synthesized_edge_ratio (0.7418), x100
\newcommand{\mathlibPin}{534cf0b}         % summary.json, snapshot.mathlib_commit; snapshot_date 2026-02-02, lean4:v4.28.0-rc1

% ---------------------------------------------------------------------------
% The dag_layer audit (Section 6).
%
% NOT IN results/. Unlike the two families flagged in Section 7's L6, these are
% recomputable in one pass over the released graph -- they ARE the check that
% section recommends -- but no CSV of them was committed. Source:
% notes/changes/2026-08-23-depth-definition-audit.md, "The inversion is
% confirmed directly against main's own data".
% ---------------------------------------------------------------------------
\newcommand{\layerZeroNodes}{119,633}   % declarations at dag_layer 0; maximum in_degree among them is 0
\newcommand{\layerCorrLogIndegree}{\ensuremath{+0.677}} % corr(dag_layer, log1p(in_degree)) = +0.6766
\newcommand{\layerTruncatedPct}{13.8}   % 42,520 of 308,129 nodes never dequeued (13.80%): 5,732 published as -1, 36,788 carrying a partial value indistinguishable from a completed one

% ---------------------------------------------------------------------------
% Limitations (Section 7).
% ---------------------------------------------------------------------------
\newcommand{\symmTransHuman}{5,222}  % results/vocabulary_contamination.csv, corpus=human, unexcluded_tactic_name_occurrences
\newcommand{\symmTransMachine}{14}   % results/vocabulary_contamination.csv, corpus=machine, unexcluded_tactic_name_occurrences

% NOT IN ANY CSV, and not regenerable: the Archive/Imo genre check was an
% ad-hoc read-only measurement, so it fails the regeneration rule by its own
% record's admission. Committed in prose only: paper/main.tex (tag
% mathai26-submitted), Limitations, and
% notes/changes/2026-09-04-archive-imo-exploratory.md. EXPLORATORY and NOT
% PRE-REGISTERED; never to be reported as a control.
\newcommand{\genreCheckRatio}{5.7}   % exploratory genre-matched ratio, Archive/Imo against the machine corpus rarefied to the same proof count

% ---------------------------------------------------------------------------
% Structural breadth (Section 5.6). Differences of committed column values,
% machine minus human, on breadth_per_step. SIGN IS STABLE AND NEGATIVE AT ALL
% THREE SCOPES on the tactic-step axis; the magnitude moves 4x, so all three
% are quoted together and none alone (CLAUDE.md section 3).
% ---------------------------------------------------------------------------
\newcommand{\breadthGapFull}{\ensuremath{-0.607}}    % 0.5139311790735444 - 1.1211735884383387; results/metrics_common_support.csv, breadth_per_step, machine/full minus human/full
\newcommand{\breadthGapSupport}{\ensuremath{-0.780}} % 0.3423192580402258 - 1.1217692880922787; results/metrics_common_support.csv, breadth_per_step, machine/support minus human/support
\newcommand{\breadthGapMatched}{\ensuremath{-0.201}} % 0.3423192580402258 - 0.543314918597603; results/metrics_common_support.csv, breadth_per_step, machine/matched minus human/matched

% Machine mass under the RETRACTED character-axis overlap region, quoted only
% to say how little of the corpus the withdrawn sign flip rested on. The flip
% itself (-0.80 -> +0.04) is a character-axis artifact and must never be
% quoted as a result: notes/changes/2026-08-19-common-support-step-axis.md.
\newcommand{\charOverlapMachineRetainedPct}{14.3} % 4,248/29,750 = 0.14279; results/metrics_by_length_slice.csv, n at machine/overlap over machine/full

% ---------------------------------------------------------------------------
% Usage concentration (Section 5.5). These are OCCURRENCE counts, not
% membership counts: a premise is counted once per use, where the coverage
% block above counts it once per corpus.
% ---------------------------------------------------------------------------
\newcommand{\humanOccurrences}{437,062}   % results/vocabulary_contamination.csv, corpus=human, occurrences_corrected
\newcommand{\machineOccurrences}{61,897}  % results/vocabulary_contamination.csv, corpus=machine, occurrences_corrected
\newcommand{\humanPremisesPerProof}{6.24}   % results/coverage_by_policy.csv, policy=corrected, human_premises_per_proof (6.2360813857261075)
\newcommand{\machinePremisesPerProof}{2.08} % results/coverage_by_policy.csv, policy=corrected, machine_premises_per_proof (2.0805714285714285)

% NOT IN ANY CSV. Per-premise occurrence shares were computed for the frozen
% MATH-AI submission and the per-premise table was not committed to results/.
% Both figures are committed in prose: paper/main.tex (tag mathai26-submitted),
% subsection "A characteristic concentration", and notes/decisions.md:1048.
% They are DESCRIPTIVE AND POST-HOC and are labelled as such where used; the
% source line records them alongside the withdrawn depth analysis, so they must
% not be read as supporting any depth claim.
\newcommand{\sqNonnegOccurrencePct}{27.0}  % share of machine premise occurrences
\newcommand{\topTwentyFiveOccurrencePct}{68.1} % top-25 machine premises, share of machine premise occurrences

% ---------------------------------------------------------------------------
% Derived figures. Each is a quotient of already-committed, already-reported
% values; the formula is printed so a reviewer can check it against the CSVs
% without reconstructing anything. No new analysis (paper-vericode/CLAUDE.md).
% ---------------------------------------------------------------------------

% Per-class extraction efficiency = achieved retention / reachable-bound retention
%   = (1 - extractor) / (1 - reachable), both columns of
%     results/instrument_penalty_decomposition.csv (miss rates, partition basis).
% Corpus-level efficiency is \efficiencyReachable = 0.8815, which sits inside
% this range.
\newcommand{\efficiencyArith}{0.8945}  % (1-0.4924816039244961)/(1-0.43265436706835875), instrument_penalty_decomposition.csv class=arith, extractor & reachable
\newcommand{\efficiencySimp}{0.8743}   % (1-0.5165661161111661)/(1-0.44703787953713775), instrument_penalty_decomposition.csv class=simp, extractor & reachable
\newcommand{\efficiencyDecide}{0.8963} % (1-0.3252788104089219)/(1-0.24721189591078066), instrument_penalty_decomposition.csv class=decide, extractor & reachable
\newcommand{\efficiencyOther}{0.8829}  % (1-0.6135211267605634)/(1-0.5622535211267605), instrument_penalty_decomposition.csv class=other, extractor & reachable
\newcommand{\efficiencyNone}{0.8871}   % (1-0.31463628396143734)/(1-0.22743207712532865), instrument_penalty_decomposition.csv class=none, extractor & reachable
\newcommand{\efficiencyByClassMin}{0.8743} % min over the five classes above (simp)
\newcommand{\efficiencyByClassMax}{0.8963} % max over the five classes above (decide)

% Rarefaction curve shape. Both corpora are still gaining vocabulary at full
% size, so both counts are lower bounds; these are the last-quartile gains,
% which say how fast each curve is flattening.
\newcommand{\rarefactionLastQuarterMachine}{103} % 1127.0 - 1024.4 = 102.6; vocabulary_rarefaction.csv, corpus=machine, vocab_mean at fraction=1.0 minus fraction=0.75
\newcommand{\rarefactionLastQuarterHuman}{10,667} % 78642.0 - 67975.4 = 10666.6; vocabulary_rarefaction.csv, corpus=human, vocab_mean at fraction=1.0 minus fraction=0.75
% Half-range of the seeded draws, as a percentage of the mean, at a quarter of
% the human corpus. NOT the widest spread on the curve: the draws are wider at
% smaller n (1.9% at fraction=0.05) and tighter at larger n (0.12% at 0.75).
% The quarter-corpus point is quoted because it is the size band the 1:1
% matched design draws in (22,785 of 70,086 human proofs, 32.5%).
\newcommand{\rarefactionDrawSpreadPct}{0.4} % (36,102 - 35,793)/2/35,910 = 0.0043; vocabulary_rarefaction.csv, corpus=human fraction=0.25, vocab_max, vocab_min, vocab_mean

% Vocabulary ratios under the two size controls.
\newcommand{\vocabRatioRarefied}{44.0} % \humanVocabRarefied / \machineVocab = 49,601/1,127 = 44.0115; vocabulary_rarefaction.csv (human, fraction=0.4244784978455041, vocab_mean) over coverage_by_policy.csv (policy=corrected, machine_vocab)
\newcommand{\vocabRatioSupport}{70.0}  % \humanVocabSupport / \machineVocabSupport = 78,591/1,123 = 69.9831; both from coverage_common_support.csv, scope=common_support
% Mathematical-layer-only basis, the robustness row for the reported
% both-layer coverage. Quotient of two committed layer counts.
\newcommand{\vocabRatioMathOnly}{89.6} % \humanVocabMath / \machineVocabMath = 75,063/838 = 89.5740; both from results/layer_composition.csv, mathematical
\newcommand{\vocabRatioMatched}{18.8}  % \humanVocabMatched / \machineVocabMatched = 21,162/1,123 = 18.8442; both from coverage_common_support.csv, scope=matched_1to1

% Fraction of each corpus retained by the common-support restriction.
\newcommand{\supportRetentionHuman}{0.9990}   % \humanProofsSupport / \humanProofs = 70,013/70,086 = 0.998958; coverage_common_support.csv, human_proofs at scope=common_support over scope=full
\newcommand{\supportRetentionMachine}{0.7659} % \matchedProofs / \machineProofs = 22,785/29,750 = 0.765882; coverage_common_support.csv, machine_proofs at scope=common_support over scope=full

% Audit sample. NOT in any CSV: results/extraction_audit.csv records only the
% 198 drawn declarations (layer=gate matcher=worst_case, total). The located
% count is committed in prose and in code:
%   notes/changes/2026-08-16-extraction-audit-first-run.md ("187 declarations
%   located, 11 not (94.4% locate rate)"), notes/decisions.md:512, and
%   src/config.py:129 (comment on AUDIT_RETENTION_RECORDED).
\newcommand{\auditDeclsLocated}{187}
\newcommand{\auditDeclsDrawn}{198}     % results/extraction_audit.csv, layer=gate matcher=worst_case, total
\newcommand{\auditTargets}{2,988}      % results/extraction_audit.csv, layer=mathematical matcher=loose, total

% ---------------------------------------------------------------------------
% Layer composition of the two premise vocabularies (M2.3, M6).
% Counts are set-membership against the frozen provenance partition
% (graph.load.split_hub_layers); no aggregation, so Li et al.'s Finding 3 does
% not bind. Both vocabularies place fully: unplaced = 0 on each side.
% ---------------------------------------------------------------------------
\newcommand{\humanVocabInfra}{3,579}   % results/layer_composition.csv, corpus=human, infrastructure
\newcommand{\humanVocabMath}{75,063}   % results/layer_composition.csv, corpus=human, mathematical
\newcommand{\humanVocabInfraPct}{4.6}  % results/layer_composition.csv, corpus=human, infrastructure_fraction (0.045510032806897076), x100
\newcommand{\humanVocabMathPct}{95.4}  % results/layer_composition.csv, corpus=human, mathematical_fraction (0.954489967193103), x100
\newcommand{\machineVocabInfra}{289}   % results/layer_composition.csv, corpus=machine, infrastructure
\newcommand{\machineVocabMath}{838}    % results/layer_composition.csv, corpus=machine, mathematical
\newcommand{\machineVocabInfraPct}{25.6} % results/layer_composition.csv, corpus=machine, infrastructure_fraction (0.25643300798580304), x100
\newcommand{\machineVocabMathPct}{74.4}  % results/layer_composition.csv, corpus=machine, mathematical_fraction (0.743566992014197), x100

% How much more infrastructure-weighted the machine vocabulary is.
\newcommand{\infraWeightRatio}{5.6}    % machine infrastructure_fraction / human infrastructure_fraction = 0.25643300798580304/0.045510032806897076 = 5.6346; both from results/layer_composition.csv, infrastructure_fraction, rows corpus=machine and corpus=human

% ---------------------------------------------------------------------------
% Residual below the reachable bound, by cause. Shares are transcribed
% percentages, not fractions, so nothing is converted; they sum to 100.0.
% ---------------------------------------------------------------------------
\newcommand{\residualTargets}{248}          % results/residual_decomposition.csv, residual_targets (mathematical targets the reachable bound finds and the extractor misses)
\newcommand{\residualResolutionPct}{89.5}   % results/residual_decomposition.csv, class=resolution_vs_substring, share_pct
\newcommand{\residualKeywordPct}{8.9}       % results/residual_decomposition.csv, class=keyword_list, share_pct
\newcommand{\residualCommentPct}{1.6}       % results/residual_decomposition.csv, class=comment_or_string_only, share_pct


\title{A Prover's Share of Mathlib: Measuring Premise Coverage
Without Running Lean}

% Shown in the title footnote, per the template's note above.
\workshoptitle{AI for Verifiable Coding @ NeurIPS 2026}

% Camera-ready: deanonymized.
\author{%
  Anirudh Balaji \\
  PES University \\
  Bengaluru, India \\
  \texttt{PES1UG25AM042@stu.pes.edu} \\
}

\begin{document}

\maketitle

\begin{abstract}
  A Lean proof the kernel accepts is correct; it is also a \emph{selection} of
  premises from a library, and \texttt{pass@k} is blind to that selection. We
  ask what can be measured about it from released artifacts alone (a
  published mathlib dependency graph and pre-verified proof corpora, with no
  Lean invocation), and how reliably. The instrument needs characterizing
  first. The release's explicit-edge flag does not mean ``written in source'',
  so no source-level extractor, however good, reaches unit retention against
  it. We therefore report retention as a scope decomposition, at \retention{}, of
  which \efficiencyReachable{} is what a per-declaration instrument can reach
  and the rest is not addressable by reading source. That efficiency is
  flat across tactic classes, so the instrument does not fail preferentially on
  the proof styles that separate the corpora. We quantify the displacement our
  own extractor induces, and show that proof length must be measured on a
  wrap-invariant axis, because character counts measure source formatting
  rather than proof size. Measured
  with that instrument, mathlib's proofs draw on \humanVocab{} distinct
  premises against \machineVocab{} for a released machine corpus:
  \vocabRatio{}$\times$ on the full corpora, \areaRestrictedRatio{}$\times$
  restricted to the library regions the machine vocabulary occupies, and
  \vocabRatioMatched{}$\times$ matched on proof count, tactic class and proof
  length, with the human-only fraction falling from \humanOnlyFullPct{}\% to
  \humanOnlyMatchedPct{}\% across those designs. Only \machineOnlyFull{}
  premises are machine-only, and \machineVocabInfraPct{}\% of the machine
  vocabulary is language infrastructure rather than mathematics. A depth
  finding we built on a released per-declaration attribute is withdrawn,
  because the attribute is a centrality proxy. The one-line check that would
  have caught it is to correlate a derived attribute against in-degree and
  inspect its extremes.
\end{abstract}

% Section 1 -- introduction, half a page. Selection framing, three
% contributions, one statement of the range. The abstract and 5.1 carry the
% rest. No literal numbers: macros only.

\section{Introduction}
\label{sec:intro}

A Lean proof the kernel accepts is correct, and nothing more needs saying about
that. But a proof is also a \emph{selection}. Of the hundreds of thousands of
declarations in a library, it reaches for a particular few. That selection
determines what a premise-selection model must learn to retrieve, what breaks
when a lemma is renamed or generalized, and whether a development participates
in the library or merely sits beside it. Correctness is settled by the kernel. Nothing settles the selection.

We make it measurable. A proof's premises are nodes in the library's dependency
graph, so a proof population induces a subgraph, and two populations can be
compared by which parts of the library their subgraphs occupy. We compare a
released machine corpus of \machineProofs{} Lean proofs against mathlib's
\humanProofs{} declarations, working entirely from published artifacts and
invoking Lean nowhere.

Three contributions follow. First, what a released-artifact pipeline can
measure and how reliably (Section~\ref{sec:instrument}): the release's
\texttt{is\_explicit} flag does not mean ``written in source'', so no
source-level instrument should be expected to reach unit retention against it;
we report retention as a scope decomposition at \retention{}, of which
\efficiencyReachable{} is what a per-declaration instrument can reach, flat
across tactic classes so the instrument does not fail preferentially on the
proof styles that separate the corpora; we quantify the displacement our own
extractor induces, and show that proof length must be measured on a
wrap-invariant axis. Second, the coverage result
(Section~\ref{sec:results}): mathlib's proofs draw on \humanVocab{} distinct
premises against the machine corpus's \machineVocab{}, a gap that survives
every control we can construct at \vocabRatio{}$\times$ on the full corpora,
\areaRestrictedRatio{}$\times$ restricted to the machine vocabulary's own
library regions, and \vocabRatioMatched{}$\times$ matched on proof count,
tactic class and length. The direction is invariant across designs and the
magnitude is not, so we quote the range throughout. Third, a caution
(Section~\ref{sec:dag-layer}): a depth finding we built on a published
attribute is withdrawn, because the attribute is a centrality proxy, and we
give the check that would have caught it.


% Section 2 -- background. Orientation for a verification reader who may not
% know mathlib's internals. Tight: everything here exists to make Sections 3-6
% readable. No literal numbers: every figure is a macro from numbers.tex.

\section{Background}
\label{sec:background}

\paragraph{mathlib4 and why its dependencies form a graph worth studying.}
Lean~4's mathematical library is a single shared development whose declarations
are each checked against a kernel, so the proof term one elaborates to names
exactly the other declarations it depends on. ``What is used to prove what'' is
recorded mechanically across the whole library, with no citation practice or
bibliographic convention standing between the record and the reasoning.

\paragraph{The published graph.} \citet{li2026network} analyse that structure at
three granularities (declarations, the modules that contain them and their
imports, and namespaces) and release the graph itself. We consume the
declaration layer and do not rebuild it: \graphDeclarations{} declarations and
\graphEdges{} dependency edges, pinned at mathlib commit \mathlibPin{}
(February 2026). Working from the release rather than from our own extraction
is what makes the human side of this paper comparable to theirs.

\paragraph{Most dependency edges were never written by anyone.} The finding that
shapes this work is that only \explicitEdgeCount{} of those edges
(\explicitEdgePct{}\%) are \emph{explicit}; the remaining
\synthesizedEdgePct{}\% are synthesized during elaboration, arriving through
instance resolution, unification and tactic machinery rather than from anything
in the source a contributor typed. A dependency graph read whole is therefore
substantially a record of what the compiler did. Li et al.\ identify the
explicit subgraph as the closer proxy for human-intended dependency structure,
and that subgraph is our object of study throughout. It also settles our
instrument. Source-visible extraction (Section~\ref{sec:extraction}) is the
correct tool for that object rather than a compromise against a fuller one.
What \texttt{explicit} does and does not mean is a separate question, and
Section~\ref{sec:is-explicit} takes it up.

\paragraph{The two corpora, and the difference between them.} Our human
reference is mathlib itself: the \humanProofs{} declarations of the pinned
checkout that carry a located source block. Our machine corpus is a released
set of \machineProofs{} Lean proofs produced by a prover we did not run
\citep{goedelprover}, and its statements come from Lean Workbook
(competition problems). The larger difference is what the two are for. A mathlib declaration is general-purpose
infrastructure written to be reused; a competition proof is a self-contained
puzzle. That difference is a property of the released corpora rather than of
provers, and it is the confound this comparison cannot remove. We state it
here rather than deferring it: Section~\ref{sec:size-controls} narrows it by
matching library area and Section~\ref{sec:limitations}, L4, bounds what the
result can mean in its presence.

\paragraph{Centrality measures technical utility, not mathematical depth.}
Their third finding constrains every structural claim we make. The most heavily
cited declarations in this graph are language infrastructure, and they identify
a two-layer hub structure separating it from mathematical results proper. We
adopt both consequences: no aggregate centrality figure without the layers
separated, and ``more central'' never read as ``less deep'', an inference
\citet{li2026network} contradict in print and
Section~\ref{sec:dag-layer} records us failing to respect.

Li et al.\ pose structural drift, as machine-generated proofs enter the library,
as future work. This paper measures a machine population against their human
baseline.


\section{The extraction instrument, characterized}
\label{sec:instrument}
We characterize the instrument before using it.
Section~\ref{sec:extraction} states what the extractor does; the four sections
after it establish what a pipeline of this kind can measure and how reliably:
what \texttt{is\_explicit} actually means
(Section~\ref{sec:is-explicit}), what retention can reach against it
(Section~\ref{sec:retention}), what displacement the extractor induces in the
comparison (Section~\ref{sec:reference-line}), and which axis proof length
must be measured on (Section~\ref{sec:length-axis}). The first two are
properties of the released artifact rather than of our code, and a pipeline
can consume either without receiving any signal that it has gone wrong.

% Section 3.1 -- what the extractor is, compressed. Mechanics (terminator set,
% tokenizer, namespace prefixes, keyword list) live in Appendix A.
% No literal numbers: every figure is a macro from numbers.tex.

\subsection{Source-visible extraction at declaration scope}
\label{sec:extraction}

The extractor takes two inputs: the mathlib checkout at the pinned revision,
and the set of declaration names in the published graph. It compiles nothing
and elaborates nothing, which is what makes it an instrument for the explicit
subgraph rather than an approximation of the full one.

Its scope is the whole declaration, statement and proof. Explicit edges come
from the elaborated term of the entire declaration, so a hypothesis contributes
premises. In \texttt{theorem foo (h : Nat.Prime p)}, \texttt{Nat.Prime} is
an explicit target. Extracting only what follows \texttt{:=} would drop
such targets while leaving them in the denominator, which is an under-scoped
instrument rather than a conservative one. The release ships no source
positions, so the block is located textually; the locate rate is measured
rather than assumed, at \auditDeclsLocated{} of \auditDeclsDrawn{} sampled
declarations.

Within that block the identifier tokens are resolved against the declaration
set, under the block-location, tokenization, resolution and prefix rules
Appendix~\ref{app:extractor} gives in full.

Two failure modes matter for reading the results. The first is deliberate: the
extractor never searches the declaration set by suffix, because resolving
\texttt{mp} that way would credit every declaration in mathlib whose name ends
in \texttt{.mp}. Tokens that resolve to nothing therefore stay out of the
premise set entirely, and the cost of that refusal is the largest component of the
shortfall below the declaration reference point of Section~\ref{sec:retention}
(\residualResolutionPct{}\%, Appendix~\ref{app:residual}). The second is
imposed: dot notation names a declaration only once the type of the local it
projects from is known, which requires elaboration, so that residue is reported
rather than resolved.


% Section 3.2 -- instrument characterization, part one.
% Companion to the dag_layer finding. Framing: a caution for consumers of the
% MathlibGraph release, not a limitation of this work. No literal numbers:
% every figure is a macro from numbers.tex. The retention decomposition lives
% in 3.3 -- this section states the set relation and points forward.

\subsection{\texttt{is\_explicit} does not mean ``typed in source''}
\label{sec:is-explicit}

The dependency graph we build on~\citep{li2026network} distinguishes \emph{explicit} edges from the
much larger set of compiler-synthesized ones, and identifies the explicit
subgraph as the closer proxy for human-intended dependency structure. That
subgraph is our object of study, and the distinction is the right one to draw.
The flag's name, however, invites a reading it does not support, and anyone
building a source-level pipeline on this release should know which reading is
correct before choosing a denominator.

\paragraph{Where the flag comes from.}
We consume the published graph rather than re-running its extraction; what
follows is provenance, read off the release's own parser and the documented
behaviour of the tool upstream of it. \texttt{is\_explicit} originates as the
\texttt{*} marker emitted by \texttt{lean-training-data}'s \texttt{premises}
command~\citep{leantrainingdata}, and the release's parser consumes it with that
meaning. The marker records that a constant occupies an \emph{explicit argument
position of the elaborated term} (a position not filled by instance
resolution or unification). It is a property of the term the kernel accepts, and
by construction not a claim about the surface syntax a human wrote, since it is
assigned after elaboration, when the source text is gone.

\paragraph{The two readings come apart at the top of the graph.}
Elaboration inserts constants of its own. In the published graph,
proof-irrelevance and rewriting plumbing (\texttt{propext}, \texttt{Eq.mpr},
\texttt{of\_eq\_true}) carry explicit in-edges throughout the explicit
subgraph and are among its most heavily cited nodes by in-degree. No mathlib
author writes any of them. They are artifacts of how \texttt{rw}, \texttt{simp}
and the equation compiler construct terms, and they are explicit in the sense
the flag means, sitting in explicit argument positions of terms that no human
typed in that form. The same mechanism operates less visibly throughout
the library, in notation-backed definitions that are never written as
identifiers, in constants emitted by arithmetic normalization, and in simp-set
lemmas that fire without being named.

\paragraph{The consequence for any source-level instrument.}
We take the narrower of two available claims, and it is enough. Some explicit
dependencies are never written as identifiers at all (the plumbing above is the
clearest case), so no instrument that reads source can recover them, and
\emph{no} source-level instrument, including a hypothetically perfect one with
exact name resolution, complete namespace and \texttt{open} handling and no
tokenization loss, should be expected to reach unit retention against the
explicit set. We do not claim the stronger set relation. The two sets are not
nested in either direction: a name written in source need not survive as the
same dependency in the elaborated term, and source text carries names that are
not cited constants at all. This is why we report retention against a scope
decomposition rather than against a threshold.
Section~\ref{sec:retention} gives that decomposition as three quantities cut
with one matcher at three scopes, and says what each does and does not
license.

Two consequences follow for anyone building on this release. High retention
rates reported for source-visible extraction elsewhere~\citep{cslibpremise} are
not comparable to ours, because a retention figure inherits the semantics of the
artifact it is scored against and belongs to an (instrument, ground-truth) pair
rather than to the instrument alone. And membership in the explicit set is not
evidence that a human wrote the name, so it should not be used to build a
retrieval target, to score a premise-selection model against ``the premises a
mathematician used,'' or to set a parser's accuracy bar. A check that costs
very little catches this, and it is the same one that catches the attribute of
Section~\ref{sec:dag-layer}: take the highest-in-degree nodes and ask whether a
mathlib author would type them.


% Section 3.3 -- retention and its scope decomposition.
% Characterization, not apology. No literal numbers: every figure is a macro
% from numbers.tex. Budget: 500-600 words (section 3 total is 2.25 pages).

\subsection{Retention, and what its denominator is a property of}
\label{sec:retention}

We report retention against a decomposition rather than against a threshold,
and the reason is instructive about how instruments in this setting should be
scored at all.

Our original design gated the project on retention: the extractor had to clear
a fixed fraction of the explicit targets in its sample or the approach was
abandoned. That gate was retired before it ever returned a verdict we acted on,
because its threshold turned out to encode two things that have nothing to do
with our parser. The first is the elaborator behaviour of
Section~\ref{sec:is-explicit}: part of the explicit set was never in the source
text, so no fraction of it is reachable by reading source. The second is
editorial: mathlib authors group related results in files, so whether a premise
name happens to appear elsewhere in a declaration's file is a fact about how the
library is organised. A threshold that moves when either of those changes is not
measuring the instrument, and a pass it awards is not evidence about the
instrument either. What replaces it is a decomposition into quantities that
are each measured the same way, so that the steps between them can be read
off rather than argued for.

\begin{table}[h]
\centering
\caption{Retention of mathematical-layer explicit targets, same sample and same
loose matcher, at three scopes. These are matcher-based reference points, not
attributable bounds: a loose name match is evidence that a name occurs within
the scope, not that the premise was used there, and it can equally miss a
premise written under another spelling or notation.}
\label{tab:retention}
\small
\begin{tabular}{llc}
\toprule
Quantity & Denominator & Value \\
\midrule
Whole-file reference   & explicit targets, name matched anywhere in the file & \ceilingWholeFile{} \\
Declaration reference  & explicit targets, name matched in the declaration   & \reachableBound{} \\
Achieved retention     & explicit targets, our extractor                     & \retention{} \\
Efficiency             & achieved against the declaration reference          & \efficiencyReachable{} \\
\bottomrule
\end{tabular}
\end{table}

Read downward, Table~\ref{tab:retention} separates three steps of very
different kinds. Targets missed even by a whole-file match are not recoverable
by reading source under any spelling the matcher knows; elaborator-inserted
constants are the mechanism we can name for that, and we do not claim it
accounts for all of them. The step from the whole-file to the declaration
reference is an organisational effect: those targets are named somewhere in the
file but nowhere in the declaration that depends on them, so a per-declaration
instrument cannot see them however good its name resolution is. Only the last
line is a claim about our parser, and it is the figure we ask to be judged on.

The \residualTargets{} targets the declaration reference point finds and the
extractor misses are identified rather than merely counted, and
Appendix~\ref{app:residual} decomposes them: the dominant class is the
definitional gap between resolving an identifier and matching a substring, and
closing it entirely would mean buying retention at the cost of precision. Two
properties make this a characterization rather than a disclosure. Each residual
class names a mechanism. And efficiency is close to flat across tactic classes,
from \efficiencyByClassMin{} to \efficiencyByClassMax{}, so the instrument does
not fail preferentially on the proof styles that distinguish the two corpora.
That is a statement about mathlib proofs only, as
Section~\ref{sec:limitations}, L1, records.


% Section 3.4 -- the instrument reference line, compressed. Diagnostics behind
% the precision term live in Appendix C. No literal numbers: macros only.

\subsection{The instrument reference line}
\label{sec:reference-line}

An extractor that recovers part of the explicit set displaces any comparison
built on its output, so the displacement has to be quantified. We do this with
a reference line, the human--machine gap the instrument produces on its own.

Two artifacts contribute and they are opposite in kind. Recall failure drops
premises a declaration really cites; precision failure admits premises it does
not. They do not cancel, because the corpora differ in the proof styles each
artifact is worst on, so both displace the comparison the same way and the line
is their sum. We therefore report it as separate terms rather than as one
averaged scalar.

The recall term is standardized jointly on tactic class and length bin, on the
step axis of Section~\ref{sec:length-axis}. It survives that control and
strengthens under it, from \recallTermClassOnly{} class-only to
\recallTermJoint{} joint, with a bootstrap interval of (\recallTermCIlow{},
\recallTermCIhigh{}) and a band of (\recallTermBandLow{},
\recallTermBandHigh{}) across the precedence orderings of the class partition.
The correction runs against our own interest, since a larger instrument
artifact raises the bar our headline result has to clear.

The precision term is the weak point. By a criterion registered before the run
it is uninformative rather than clean. We therefore carry it as an open item
rather than folding it into the line, and Appendix~\ref{app:reference} gives
the diagnostics behind that decision.

Two limits on what this line licenses. Per-class efficiency is close to flat
(Section~\ref{sec:retention}), which is a stronger neutrality statement than
the raw per-class miss rates suggest, but it is measured on human proofs,
and nothing here establishes comparable recovery on machine proofs, which
Section~\ref{sec:limitations}, L1, states as the paper's load-bearing
assumption. And the line matters unevenly: the coverage result of
Section~\ref{sec:coverage} sits far outside the regime these artifacts operate
in, while the secondary characterizations sit inside it.


% Section 3.5 -- the length axis, compressed. The tactic-step proxy and its
% under-recording are in Appendix D. Framing: a reusable methods result, not a
% caveat about our data. No literal numbers: macros only.

\subsection{Length must be measured on a wrap-invariant axis}
\label{sec:length-axis}

Controlling for proof length is mandatory here. Short proofs mechanically cite
fewer and more common premises, so an uncontrolled length difference can
manufacture every effect we report. The choice of length \emph{axis} is itself
a measurement question with a checkable answer, and that is the reusable part
of this section.

Two axes are available. On characters the corpora barely overlap, at an
overlap coefficient of \charOverlapSample{}. On tactic steps they nearly
coincide, at \stepOverlapCensus{}.\footnote{The two rest on different bases:
\charOverlapSample{} is computed on the standardization sample the joint
control was cut on, \stepOverlapCensus{} over the full census of both corpora.
The step-axis figure on the sample basis is \stepOverlapSample{}, so the basis
is not what separates the axes.} The disagreement is explained by how each
corpus reaches the page. Mathlib source is hand-wrapped near a column limit and
carries docstrings; a released machine proof is generated in one pass, wrapped
differently or not at all, and here carries a large volume of natural-language
commentary. Character counts absorb all of that; tactic counts do not, because
wrapping and commentary do not change how many tactics run.

Standardizing on characters is therefore substantially standardizing on source
formatting, and that is not a neutral error. On the character axis one of our
secondary metrics changed sign and the recall term of
Section~\ref{sec:reference-line} collapsed to a null, and both reverse on the
step axis. Every length control in this paper accordingly runs on tactic steps,
and character length appears once, here, as the diagnostic that exposed the
problem. It is not an alternative specification, because a specification that
measures wrapping is not a second opinion about proof size. Counting steps
needs its own instrument and the obvious one under-records, which
Appendix~\ref{app:steps} measures and Section~\ref{sec:limitations}, L2,
carries as a limitation.

The check that settles the axis is cheap. Compute the overlap coefficient on
each candidate axis before choosing one. Where the axes disagree as sharply as
they do here, the choice is load-bearing, and making it after seeing the
comparison is a researcher degree of freedom.


% Section 4 -- formal definitions. Definitional only: no figures, no macros.
% The "this is set cardinality, not structure" objection is raised here and
% answered here, which is cheaper than defending it after a reviewer raises it.

\section{Structural coverage: formal definitions}
\label{sec:definitions}

Let $G$ be the explicit subgraph of the mathlib dependency graph, as
characterized in Section~\ref{sec:is-explicit}: its nodes are declarations and
its edges are the source-attributable dependencies between them. A proof $p$
selects a set of premises $P(p) \subseteq V(G)$, the declarations it cites, and
those premises span an induced subgraph $G[P(p)]$. For a corpus $C$ we write
$V(C) = \bigcup_{p \in C} P(p)$ for its \emph{vocabulary}: the node set of the
subgraph the corpus's whole proof population induces.

\paragraph{Coverage.} Given a human corpus $H$ and a machine corpus $M$,
coverage is the family of quantities over $V(H)$ and $V(M)$: the two
vocabulary sizes and their ratio; the Jaccard index
$J = |V(H) \cap V(M)| \,/\, |V(H) \cup V(M)|$; and the human-only fraction
$|V(H) \setminus V(M)| \,/\, |V(H)|$. These quantities read no edges.
Computing $J$ consults only membership, so coverage is the node-set family and
breadth, below, the edge-derived one. Appendix~\ref{app:construct} states what
the graph does and does not supply to the coverage arithmetic, why we chose a
measure that needs no aggregation, and how each quantity is reported.

\paragraph{Breadth.} Structural breadth is the topological complement: the
number of distinct communities of a fixed partition of $G$ that a premise set
touches, optionally per unit of proof length. This does read edges, since the
partition is computed from them.

\paragraph{Which definitions need the hub-layer split.} Coverage does not. It
aggregates nothing, so Finding~3 does not bind, and the layer composition of
each vocabulary is reported in Section~\ref{sec:layer-composition} as a result
in its own right rather than as a correction. Breadth and every
centrality-flavoured quantity do need it. Breadth carries a second caveat, the
partition skew of Appendix~\ref{app:breadth}, and it is reported as secondary
throughout, with its three design-dependent magnitudes quoted together, never
one alone.


\section{Results}
\label{sec:results}
This section applies the definitions of Section~\ref{sec:definitions} to the
two corpora of Section~\ref{sec:background}: coverage and the two
controls that answer the objections to it
(Sections~\ref{sec:coverage} and \ref{sec:size-controls}), then
breadth (Section~\ref{sec:breadth}), reported as unstable. The composition and
concentration of the machine vocabulary are reported in
Appendices~\ref{sec:layer-composition} and \ref{sec:concentration}. Two conventions
hold throughout. Every figure is on the corrected premise sets. Two artifacts
that resolved on the human side only (a declaration's own name and the
\texttt{lemma} keyword) are removed before any comparison, at a large cost
to the human count, and Section~\ref{sec:limitations}, L2, gives the
accounting. 

% Section 5.1 -- vocabulary coverage. The headline result.
% Ported from the frozen MATH-AI results section and expanded: intersection
% counts alongside every Jaccard so each row is recomputable, and the fifth
% design (common support) that four pages had no room for.
% No literal numbers: every figure is a macro from numbers.tex.
% The uncorrected human count is retracted and has no macro by design.

\subsection{Vocabulary coverage}
\label{sec:coverage}

Over the full corpora the two vocabularies intersect in \intersectionFull{}
premises, giving a Jaccard index of \jaccardFull{}. \humanOnlyFullPct{}\% of
the premises mathlib's own proofs cite are never cited by any machine proof in
the corpus. The asymmetry
is close to containment. Only \machineOnlyFull{} premises are machine-only, so
the machine vocabulary is very nearly a subset of the human one rather than a
different region of the library. Two properties of that small vocabulary are
reported in the appendix. It is weighted toward language infrastructure rather
than mathematics, at \machineVocabInfraPct{}\% against the human corpus's
\humanVocabInfraPct{}\%, leaving a mathematical working vocabulary of
\machineVocabMath{} declarations (Appendix~\ref{sec:layer-composition}). And
it is used very unevenly, with proofs closing in at most two steps under an
arithmetic tactic set making up \degenerateMachineSharePct{}\% of the machine
corpus against \degenerateHumanSharePct{}\% of mathlib
(Appendix~\ref{sec:concentration}). Every uncorrected instrument asymmetry we
have identified narrows the gap reported here rather than widening it, so these
figures are floors; Appendix~\ref{app:limitations}, L2, enumerates them.

The designs differ in what they are allowed to explain away. The full corpora
control for nothing. Rarefaction removes the corpus-size advantage by drawing
the human corpus down to the machine corpus's proof count. Common support
discards the cells of the (tactic class $\times$ tactic step) grid that only
one corpus occupies, so no rate is imputed where the other side has no
population. The area-restricted design keeps only the regions of the library
the machine vocabulary occupies at all, removing the objection that the two
corpora are simply doing different mathematics. The matched design is the
strictest: equal proof counts, equal tactic class, equal step bin.

\begin{table}[t]
\centering
\small
\begin{tabular}{lrrrrrr}
\toprule
design & human & machine & ratio & Jaccard & $|H \cap M|$ & human-only \\
\midrule
full corpora & \humanVocab{} & \machineVocab{} & \vocabRatio{}$\times$ & \jaccardFull{} & \intersectionFull{} & \humanOnlyFullPct{}\% \\
\textbf{area-restricted} & \textbf{\humanVocabAreaRestricted{}} & \textbf{\machineVocab{}} & $\mathbf{\areaRestrictedRatio\times}$ & \textbf{\areaRestrictedJaccard{}} & \textbf{\intersectionAreaRestricted{}} & \textbf{\areaRestrictedHumanOnlyPct{}\%} \\
1:1 class- and step-matched & \humanVocabMatched{} & \machineVocabMatched{} & \vocabRatioMatched{}$\times$ & \jaccardMatched{} & \intersectionMatched{} & \humanOnlyMatchedPct{}\% \\
\bottomrule
\end{tabular}
\caption{Premise vocabulary under three designs of increasing severity. Each
row carries its intersection count, so its Jaccard index is recomputable from
the row alone as $|H \cap M| / (|H| + |M| - |H \cap M|)$. The area-restricted
row is the primary control (Section~\ref{sec:size-controls}); the matched row
is the strictest design we ran (Section~\ref{sec:size-controls}). Two further
designs are reported in the text rather than here: rarefying the human corpus
to the machine corpus's size gives \vocabRatioRarefied{}$\times$
(Section~\ref{sec:size-controls}), and restricting to common support gives
\vocabRatioSupport{}$\times$ at a Jaccard of \jaccardSupport{}, against
\jaccardFull{} on the full corpora.}
\label{tab:coverage}
\end{table}



The point of the table is the invariance rather than any single row. The ratio
never leaves the range \vocabRatioMatched{}$\times$ to \vocabRatio{}$\times$,
and the human-only fraction never falls below \humanOnlyMatchedPct{}\%, under
controls that variously remove corpus size, proof length, tactic style and
subject area. The magnitude is design-dependent and we quote the range rather
than the largest figure; the direction and order of magnitude are not.
Coverage is also the most stable quantity we measure. Moving from the full
corpora to common support changes the Jaccard index only in its fourth decimal
place, \jaccardFull{} to \jaccardSupport{}, while that restriction discards
roughly a quarter of the machine corpus. The result is therefore not an
artifact of which proofs happen to be comparable. Every proof on both sides of
this comparison is verified. Correctness is satisfied throughout, and none of
the structure in this table is visible to \texttt{pass@k}.


% Sections 5.2 and 5.3 merged: the two controls that answer the two objections
% to the headline, stated with their results rather than with their
% construction. Design details live in 5.1's table caption and the appendix.

\subsection{Corpus size and region selection}
\label{sec:size-controls}

Two objections do not require argument, because both quantities can be varied
directly. The first is arithmetic: more proofs find more distinct premises
whatever the populations are doing, so a vocabulary compared at unequal $n$
confounds breadth with sample size. We rarefy the human corpus (seeded draws
of $n$ proofs, distinct premises counted, mean reported), which gives its
vocabulary at every corpus size. At the machine corpus's own size,
\machineProofs{} proofs, the human vocabulary is \humanVocabRarefied{}
premises, so equalizing size moves the ratio from \vocabRatio{}$\times$ to
\vocabRatioRarefied{}$\times$ rather than to one. Further down the curve the
point is sharper still. At a twentieth of the human corpus,
\rarefiedFivePctProofs{} proofs and far fewer than the machine corpus has, the
human vocabulary is \humanVocabRarefiedFivePct{} premises, an order of
magnitude above what the machine corpus reaches at full size. Neither curve has
saturated, so both reported vocabularies are lower bounds on their
populations', and the machine curve is the nearer of the two to flattening; we
do not extrapolate either beyond the corpora we have.

The second objection is the genre difference named in
Section~\ref{sec:background}: if competition problems occupy a small corner of
the library, a narrow machine vocabulary is a fact about which problems were
posed. The control we would prefer, a same-problem comparison, does not exist.
The only human Lean~4 proofs of miniF2F statements are a handful released as a
prover's own few-shot prompt pool, the LeanDojo Lean~4 port carries
\texttt{sorry} throughout, and the original corpus is Lean~3. We therefore
restrict by region instead, keeping every top-level region that contains at
least one machine premise. The rule was fixed before the run, is support-based so there
is no threshold or top-$k$ to tune, and introduces no metric. The machine
vocabulary occupies \areaRegionsMachine{} of the \areaRegionsHuman{} regions
the human vocabulary spans, and those regions hold \areaHumanSurvivalPct{}\% of
human premises, so the restriction discards most of the human vocabulary rather
than trimming it. Inside them the ratio is \areaRestrictedRatio{}$\times$ at a
Jaccard of \areaRestrictedJaccard{}, with \areaRestrictedHumanOnlyPct{}\% of
the restricted human vocabulary still untouched.

Machine leakage is zero by construction under this rule, which is what makes
it primary and its registered sensitivities the check on it. The
top-twelve-namespace variant gives \areaTopTwelveRatio{}$\times$ but discards
\areaTopTwelveMachineLeakPct{}\% of the machine vocabulary on the way, and a
restriction that drops nearly a third of the population it controls for is
answering a different question. We report the range
\areaTopTwelveRatio{}$\times$ to \areaRestrictedRatio{}$\times$ rather than
either endpoint. Neither control reaches genre, or selection \emph{within} a
region. \texttt{Nat} holds competition arithmetic and analytic number theory
alike. The closest genre check we have is exploratory, and we state it here rather than
only in limitations because it is the design that most constrains what the
headline can mean. Mathlib's own IMO archive, against the machine corpus
rarefied to the same proof count, gives \genreCheckRatio{}$\times$, the
smallest gap of any design we ran. It was not pre-registered, it is a lower
bound, it is not extraction-audited, and it is not regenerable from the
released artifact (Section~\ref{sec:limitations}, L6). These designs bound how much of the gap corpus size and region selection
explain; Section~\ref{sec:limitations}, L4, bounds what the result can mean in
the presence of what they do not.


% Section 5.6 -- structural breadth, one paragraph. Direction in the main text,
% magnitudes and masses in Appendix E. Reported as unstable, NOT as evidence.

\subsection{Structural breadth}
\label{sec:breadth}

Breadth, as defined in Section~\ref{sec:definitions}, counts the distinct
communities of the partition a proof's premises touch, per tactic step. Machine
proofs touch fewer communities per step than human ones at every scope we can
measure, and the direction is the only part of that we report: the gap is
\breadthGapFull{} on the full corpora, \breadthGapSupport{} under common
support and \breadthGapMatched{} matched 1:1, the last two computed on
\supportRetentionMachine{} of the machine corpus. The magnitude therefore moves
by a factor of four across designs that move the coverage result in its fourth
decimal place, and an earlier sign flip was an artifact of a region cut on
character length, which Section~\ref{sec:length-axis} rejects as a formatting
control. We report the metric as unstable, having defined it before we looked,
and we do not offer it as evidence for or against the coverage result, which
stands on its own controls. Appendix~\ref{app:breadth} carries the masses, the
excluded region and the partition skew that keep it secondary.


% Section 6 -- the withdrawn finding, compressed to a negative-result
% paragraph. Framing kept: the error was OURS. The concession that the
% attribute may behave exactly as documented is now a pointer to Appendix I
% (app:daglayer), not stated here. Never quote the withdrawn depth figures.

\section{A withdrawn finding: \texttt{dag\_layer} and the catching check}
\label{sec:dag-layer}

The graph release publishes a per-declaration attribute, \texttt{dag\_layer},
which we took to be distance from primitives. We built a depth metric on it,
measured a machine--human gap, and promoted that gap to a stated
counter-finding. The column is not a depth measure as we used it: layer zero
holds \layerZeroNodes{} declarations whose maximum in-degree is zero (uncited
leaves, not axioms), the highest layer holds \texttt{Eq.refl}, \texttt{Set}
and \texttt{LE.le}, and the column correlates \layerCorrLogIndegree{} with log
in-degree, so a high value means heavily cited infrastructure and the scale runs
opposite to the direction we read into it. The error was ours, and
Appendix~\ref{app:daglayer} gives the reading of the column under which it
behaves as documented, together with a second hazard for other consumers of the
release. What we measured is therefore a centrality proxy, and a
difference in it is a statement about citation weight rather than derivation
length. That is the inference \citet{li2026network} rule out in their
Finding~3, which we cite in Section~\ref{sec:background} and then walked into
on a different column of the same release. We withdraw the metric rather than
re-pointing it at a claim about centrality that it would support, because
choosing a quantity's meaning after seeing its comparison is a researcher
degree of freedom. The check that would have caught it is cheap enough to run
on any released structural attribute before building on it: correlate the
attribute against in-degree, then inspect its extremes and ask whether the
declarations at each end match the meaning its name implies.


% Section 7 -- Limitations. Ported from paper/limitations.md (L1-L6) with the
% material the four-page limit cut restored. Ordered by how much damage each
% objection does if we do not raise it first. No literal numbers: every figure
% is a macro from numbers.tex.

\section{Limitations}
\label{sec:limitations}

Three objections are load-bearing, and we state them first:
the corpora target different mathematics (L4), the instrument correction is
measured on one side and transferred to the other (L1), and everything rests on
one machine corpus (L6).

\paragraph{L1. The correction transfers a human rate to a larger population.}
Machine-side miss rates are imputed rather than measured, from per-class rates
measured on human proofs and standardized onto the machine tactic mix.
Arithmetic-tactic proofs are \machineArithMixPct{}\% of the machine corpus and
\humanArithMixPct{}\% of mathlib, so most of the machine corpus is standardized
against a stratum that is a fraction of a percent of mathlib. That rate is
near-censused and precise; the exposure is that the two populations'
\texttt{nlinarith} may not be the same object. In mathlib it is typically one
step closing a residual goal; in the machine corpus it is frequently the whole
proof. If the extractor's blind spot scales with how much of a proof
the automation closes rather than with its presence, the transfer
under-corrects, leaving machine proofs looking narrower than they are, which
biases toward our own hypothesis. Flat efficiency across human tactic classes
(Section~\ref{sec:reference-line}) constrains the instrument's behaviour on
mathlib proofs only; nothing measured here establishes comparable recovery on
machine proofs, whose dominant stratum has almost no human counterpart. We
therefore treat the correction as a lower bound and require the gap to clear a
multiple of it, not merely exceed it. Elaboration-level extraction on the machine corpus would settle it. We did not
do it because it means reproducing the publishers' extraction pipeline on a
second corpus (building the premise tool against the same toolchain and
wrapping each standalone proof as a compilable module), and because those
proofs were verified against their releasers' mathlib rather than our pinned
commit, so a fraction would fail on renames and deprecations, giving version
skew rather than a miss rate.

\paragraph{Instrument limitations.} Four further entries are stated in
Appendix~\ref{app:limitations}: L2, the asymmetries left in the accounting and
why every one of them narrows the reported gap; L3, how much of the explicit set source-visible extraction reaches; L3d,
that Table~\ref{tab:retention}'s quantities are matcher-based; and L3e, that
our reading of the flag is not validated against elaborated terms.

\paragraph{L4. The corpora target different mathematics.} Machine corpora
target competition statements and mathlib targets library development, so part
of L1's tactic-mix difference belongs to the problem sets, not the provers.
Section~\ref{sec:size-controls} narrows this by matching library
\emph{area} without removing it, because area is not \emph{genre}. A mathlib
theorem is general-purpose infrastructure; a competition problem is a
self-contained puzzle. As an exploratory, not pre-registered check, mathlib's own IMO archive gives
\genreCheckRatio{}$\times$, the smallest gap of any design we ran. It is a
lower bound (most of those declarations cite local helper lemmas outside the
graph's build target), it is not extraction-audited, and it is not
regenerable from the released artifact (L6). We report it because we measured
it, and publishing only the designs that give larger gaps would be the
objectionable choice. We therefore claim a difference
between these two released corpora under the restrictions named, not a general
property of how provers use mathlib when proving comparable theorems.

\paragraph{L6. Scope, and what is not regenerable.} One machine corpus,
\machineProofs{} proofs. These are comparisons of two distributions and not
causal claims. We write ``machine proofs concentrate on'', never ``provers are
biased toward''. Two families of figures carried from the prior version are \emph{not}
reproducible from the released artifact: the per-premise occurrence shares of
Section~\ref{sec:concentration} and the genre-check figures in L4. Both were
computed ad hoc, neither table was
committed, and no released script regenerates them. Every other number comes from a committed result file, bar
Section~\ref{sec:dag-layer}'s graph-audit figures, which any reader recomputes
from the released graph in one pass. We flag them rather than recompute them,
because recomputing after the fact is how a frozen scope stops being one.


% Section 8 -- related work, half a page. Every citation kept with its
% differentiating claim; the paper-by-paper elaboration is gone.

\section{Related work}
\label{sec:related}

\citet{li2026network} give the multilayer analysis of mathlib4 this work
consumes (the graph, the explicit/synthesized distinction, the two-layer hub
structure) and pose structural drift from machine-generated proofs as future
work; we extend their human baseline to a machine population rather than building a
baseline.

Premise retrieval asks which declarations should close a given goal and scores
a ranking against the premises a human cited. \citet{leandojo} supply that
setting and its annotations, and \citet{cslibpremise} establish source-visible
extraction for it, the precedent for our instrument. We run the relation
backwards, asking what the selection of already-verified proofs looks like in
aggregate against a library, with no right answer per goal. A retriever can
score well on every goal it is given while the population it serves still
occupies a small region of the library.

Three lines measure past correctness and each differs in its unit.
\citet{theorembench} score partial progress \emph{within} a theorem;
\citet{proofrank} characterize a proof text through proxies and pairwise
comparison, needing a preference model where ours needs none;
\citet{verisoftbench} correlate success with transitive dependency closure
size, a magnitude where coverage is a location, in a domain where the graph is
a build artifact. Our corpora are the released \citet{goedelprover} proof set
and mathlib itself; \citet{minif2f} appears for what it lacks, the human Lean~4
proofs that would permit a same-problem comparison.


% Section 9 -- reproducibility and disclosure, one paragraph.
% Double-blind: the artifact is described, never located.

\section{Reproducibility and disclosure}
\label{sec:repro}

The anonymized supplement contains the extraction and analysis source, the
committed result files behind every table, the fetch scripts that obtain the
input data we cannot redistribute, and the test suite in which the extraction
audit asserts retention against its recorded value within a stated tolerance,
so a change to the extractor fails the suite instead of silently moving a
number here. Those scripts pull the published graph and the machine corpus from
their public releases at pinned revisions and pin the mathlib checkout to the
commit the graph was built from, every analysis runs from those pins and one
fixed seed, and reported figures fall into three provenance classes that
Section~\ref{sec:limitations}, L6, names and distinguishes. LLM assistance was
used across the categories the call asks about: implementing the extraction
pipeline, the metric definitions and the analysis scripts; data processing;
running experiments; and drafting this paper; no number reported here is quoted
from a model, and the author takes responsibility for the content, including
the withdrawn finding of Section~\ref{sec:dag-layer}.


\begin{ack}
This work received no external funding.
\end{ack}

% Small references, as in the frozen submission. Grouped so it does not
% also shrink the checklist that follows.
{\small
\bibliographystyle{plainnat}
\bibliography{refs}
}

% Technical appendix. Mechanics and diagnostics moved out of the main text so
% Section 3 states what the instrument establishes rather than how it is built.
% Unlimited appendix pages are permitted; the main-text limit is nine.
% No literal numbers: every figure is a macro from numbers.tex.

\appendix

\section{Extractor mechanics}
\label{app:extractor}

\paragraph{Locating a declaration block.} The release ships no source
positions, so a declaration's block is found textually: the declaration
keyword under any attribute or modifier prefix, extending to the next line that
begins a top-level command. The terminator set is deliberately broad, because
the two failure modes are not symmetric. Over-matching truncates a block and
loses tokens at its tail, while under-matching swallows the following
declaration and invents them. Failures to locate are dominated by auto-named
instances, which have no textual declaration site.

\paragraph{Tokenization.} Comments and string literals are blanked to spaces,
preserving token boundaries; Lean's block comments nest, so they are scanned
rather than matched by regular expression. Identifier tokens are then taken in
Lean's own shape, dot-composed, with a leading dot excluded. The \texttt{.mp}
in \texttt{h.mp} is projection on a local, not a reference to a root-level
\texttt{mp}. A short keyword list drops tactic names before
resolution. It is a precision aid rather than the safety net. Every surviving
token must still resolve against the declaration set, and the list is
deliberately short, since an over-long one silently excludes real premises that share a name
with a tactic. Section~\ref{sec:limitations}, L3b, reports the one inconsistency
we left in it rather than extending it after seeing which names flattered the
result.

\paragraph{Resolution.} The rule is three steps, first hit wins: the token is a
fully-qualified declaration name; or it completes to exactly one declaration
under a namespace prefix in scope; or it resolves to nothing and is kept aside
as a bare tail. Bare tails never enter the premise set, for the reason
Section~\ref{sec:extraction} gives.

\paragraph{Namespace prefixes.} The prefixes a bare token may complete under
are the proper prefixes of the declaration's own name, together with the
\texttt{open} clauses appearing in its file. Lean scopes \texttt{open} to
sections and this does not, so the rule is over-permissive, and it can
complete a token under a namespace not actually open at that point. The cost
is precision and never recall, and it is carried in the reference line of
Section~\ref{sec:reference-line}.

\section{Residual below the declaration reference}
\label{app:residual}

The \residualTargets{} mathematical targets that the declaration reference
point finds and the extractor misses decompose into three classes.
\residualResolutionPct{}\% is a definitional gap: the reference point
substring-matches premise names while the extractor resolves identifiers, and a
resolved identifier is a stricter object than a matched substring, so closing
that gap entirely would mean abandoning resolution and buying retention at the
cost of precision. \residualKeywordPct{}\% is the over-broad keyword list of
Appendix~\ref{app:extractor}, which suppresses names that are simultaneously
tactic names and library declarations; we report its cost rather than tuning
the list after seeing which names it flatters. The remaining
\residualCommentPct{}\% occurs only inside comments or string literals, which
the extractor strips by design. The classes are identified rather than merely
counted, so each is a statement about a mechanism.

Per-class extraction efficiency (achieved retention against the declaration
reference point, within each tactic class) ranges from
\efficiencyByClassMin{} to \efficiencyByClassMax{}, with the corpus-level
figure \efficiencyReachable{} inside that range. Section~\ref{sec:limitations},
L1, states what this does and does not license about the machine corpus.

\section{Instrument reference line: diagnostics}
\label{app:reference}

The precision term is carried as an open item in
Section~\ref{sec:reference-line}, and the diagnostics behind that decision are
these. Length-controlled on the tactic-step axis the term falls to
\spuriousStepAxis{}, against \spuriousCharAxis{} on the character axis, a
near-collapse that read alone would be good news about the instrument. By a
criterion registered before the run, a control whose collapsed cells carry a
large share of the machine mix is uninformative rather than clean: here they
carry \spuriousCollapsedMachineMixPct{}\% of it, dominated by the arithmetic
stratum of Section~\ref{sec:concentration}. The character axis fails the
companion criterion of near-disjoint supports, its overlap coefficient being
\charOverlapSample{}. One diagnostic points somewhere real. Across the two
axes the imputed human rate is stationary while only the machine rate moves,
which is the signature of a compositional correction rather than of two
extrapolations converging, and that makes the step figure the more credible of
the two without making it something to freeze on.

\section{The tactic-step proxy}
\label{app:steps}

Counting tactic steps needs an instrument and the obvious one under-records.
The published tactic lists capture only main-path tactics; everything after
\texttt{<;>}, inside \texttt{induction \dots\ with} case branches, and in
nested \texttt{(by simp \dots)} term positions is missing. That is measured
rather than assumed. Against a string matcher over the same declarations the
disagreement is almost entirely one-directional, the matcher finding tactics
the lists omit and rarely the reverse. Both corpora are therefore counted with
a single fixed tactic vocabulary and one matcher. It is a size proxy, not a
taxonomy, and a tactic missing from the vocabulary costs sensitivity
symmetrically. Section~\ref{sec:limitations}, L2, states the consequence for
every stratum and mix in the paper.

\section{Structural breadth by design}
\label{app:breadth}

\begin{table}[h]
\centering
\caption{Breadth per tactic step, machine minus human, at three scopes, with
the share of the machine corpus each design retains. The sign does not move;
the magnitude moves by a factor of four, which is why
Section~\ref{sec:breadth} quotes a direction and not a magnitude.}
\label{tab:breadth}
\small
\begin{tabular}{lcc}
\toprule
Design & Gap & Machine corpus retained \\
\midrule
Full corpora   & \breadthGapFull{}    & all \\
Common support & \breadthGapSupport{} & \supportRetentionMachine{} \\
Matched 1:1    & \breadthGapMatched{} & \supportRetentionMachine{} \\
\bottomrule
\end{tabular}
\end{table}

Two things keep breadth secondary. The community partition is heavily skewed,
with most communities tiny and a few holding much of the graph, so the unit is
coarse before any comparison is made; and a quantity whose magnitude depends this strongly on the
design cannot carry weight on its own. We report it nonetheless because the
metric was defined before we looked, and omitting an unstable result would be
selection rather than caution. Breadth also measures how far apart a
population's premises lie within the library, rather than how much of the
library that population reaches, so it is not a second reading of coverage.

The excluded region under the last two designs is the set of cells with no
human counterpart, overwhelmingly the arithmetic stratum of
Section~\ref{sec:concentration}. An earlier version of this comparison did
move in sign. Restricting to a length-overlap region gave a positive gap
against the full corpora's negative one. That region was cut on character
length, which
Section~\ref{sec:length-axis} shows to be a formatting control rather than a
length control, and it retained \charOverlapMachineRetainedPct{}\% of the
machine corpus. On the wrap-invariant axis the sign is stable at all three
scopes, so the flip is an artifact of the rejected axis and not a finding.

\section{Limitations L2, L3, L3d and L3e}
\label{app:limitations}

\paragraph{L2. Instrument asymmetries, and their direction.} Human tactic use
comes from the publishers' extracted lists and machine tactic use has no
equivalent, so both sides are cut with the string matcher of
Section~\ref{sec:length-axis} and Appendix~\ref{app:steps}; the two instruments
disagree materially and in one direction, and one false positive is kept rather
than removed after the fact, because the tactic list was frozen before the
machine corpus was touched. On the extraction side, two artifacts that
resolved on the human side only (a declaration's own name and the
\texttt{lemma} keyword) are removed before any reported figure, at a large
cost to our own headline, and three further asymmetries are left in place: \texttt{symm} and
\texttt{trans} resolve \symmTransHuman{} times on the human side against
\symmTransMachine{}, prefix completion supplies far more premises per proof to
the human side because every machine proof carries the same fixed namespace
header, and coverage is counted over both hub layers where the mathematical
layer alone would give \vocabRatioMathOnly{}$\times$ rather than
\vocabRatio{}$\times$. Every uncorrected asymmetry we have found narrows the
reported gap, so Section~\ref{sec:coverage}'s figures are floors. The
corrections we applied cost us, and the ones we declined would have paid us.
Separately, no aggregate centrality figure appears without the hub layers
separated and we never infer shallowness from centrality;
Section~\ref{sec:dag-layer} reports where we broke that rule ourselves.

\paragraph{L3. Extraction cannot recover the explicit set in full.} Sections
\ref{sec:is-explicit} and \ref{sec:retention} characterize this rather than
bound it away. The whole-file reference point is \ceilingWholeFile{},
which no per-declaration instrument beats, and it is a loose-matcher figure
rather than an attributable bound. The reported vocabularies are
instrument-filtered.

\paragraph{L3b/L3c. The asymmetries that remain all shrink the gap.} Two
artifacts that resolved on the human side only, a declaration's own name
and the \texttt{lemma} keyword, are removed before any reported figure, at a
large cost to our own headline. Three further asymmetries are
left in place and all run the same way. \texttt{symm} and \texttt{trans} are
tactic names that are also mathlib theorems, resolving \symmTransHuman{} times
on the human side against
\symmTransMachine{}; we leave
them rather than extend a keyword list after seeing which names flatter the
result. Prefix completion supplies far more premises per proof to the human
side,
since every machine proof carries the same fixed namespace header while a
mathlib declaration sits in its own. That is legitimate resolution rather than
contamination. And coverage is counted over both hub layers, where the
mathematical layer alone would give \vocabRatioMathOnly{}$\times$ rather than
\vocabRatio{}$\times$. Every uncorrected asymmetry we have found narrows the
reported gap, so Section~\ref{sec:coverage}'s figures are floors. The
corrections we applied cost us, and the ones we declined would have paid us.

\paragraph{L3d. The reference points are matcher-based.} The quantities of
Table~\ref{tab:retention} are cut with a loose name matcher: a match is
evidence that a name occurs within the scope, not that the premise was used
there, and we do not measure how often it credits an incidental occurrence.

\paragraph{L3e. Our reading of the flag is not validated against elaborated
terms.} The semantics established in Section~\ref{sec:is-explicit} come from
the release's own parser and the documented behaviour of the tool upstream of
it, not from elaborating proofs ourselves, which we do not do anywhere in this
work.

\section{A quarter of the machine vocabulary is not mathematics}
\label{sec:layer-composition}

The retention figures of Sections~\ref{sec:retention} and
\ref{sec:reference-line} are quoted on the mathematical layer, using a
partition of the declaration set by provenance rather than by namespace
prefix: Lean core constants, instances, coercions, kernel-level kinds and
compiler-generated names are infrastructure, and the rest is mathematical
content. Provenance is what makes the partition usable. A prefix rule misses
every root-level elaborator constant, and it over-captures genuine mathlib
results whose names happen to sit under an infrastructure namespace.
Applied to the premise vocabularies themselves rather than to retention
targets, the same partition says something about the corpora.

The human vocabulary is \humanVocabInfraPct{}\% infrastructure:
\humanVocabInfra{} declarations against \humanVocabMath{} mathematical ones.
The machine vocabulary is \machineVocabInfraPct{}\% infrastructure,
\machineVocabInfra{} against \machineVocabMath{}, a factor of
\infraWeightRatio{} more infrastructure-weighted. Both vocabularies place
completely, with no declaration left unclassified on either side.

The machine vocabulary therefore differs from the human one in composition,
not only in size. A quarter of it is the plumbing of
Section~\ref{sec:is-explicit} (proof-irrelevance and rewriting machinery that
every Lean proof drags in and no author selects), so the mathematical working
vocabulary of the machine
corpus is \machineVocabMath{} declarations. That is the figure to hold in mind
next to mathlib's \humanVocabMath{}, since it is the set of mathematical
results the released corpus actually reaches for, as opposed to the set it
cites.

This is a membership count and nothing more. Each premise is classified once,
by provenance, and the two sides are compared by how many fall in each class.
There is no aggregation over a centrality score, so Li et al.'s Finding~3 does
not bind here and we draw no inference about mathematical depth from it. A
vocabulary weighted toward infrastructure is not thereby a shallower one.

This composition sits next to the concentration reported in
Section~\ref{sec:concentration}, where most of the machine corpus is short
automation calls. A proof that discharges its goal in one such call still
elaborates into a term, and that term carries the rewriting and
proof-irrelevance machinery with it, so the two measurements are consistent
with one picture of the corpus. We report them as two measurements rather than
as one explained by the other. For evaluation the consequence is direct. A
vocabulary that is \machineVocabInfraPct{}\% infrastructure and one that is
\humanVocabInfraPct{}\% are different objects, and only the second is evidence
of mathematical range, a distinction no correctness-based metric records.

The split also settles a basis question left open in
Section~\ref{sec:coverage}. Coverage is reported over both layers. Counting it
on the mathematical layer alone would give \vocabRatioMathOnly{}$\times$
rather than \vocabRatio{}$\times$, with a correspondingly higher human-only
fraction, because the machine vocabulary is more infrastructure-weighted and
so loses proportionally more under the restriction. We report the both-layer
basis, which is the weaker of the two for our own hypothesis, and print the
other here. Where two bases are defensible, the one to keep is the one that
raises the bar the result has to clear.

\section{A characteristic concentration}
\label{sec:concentration}

Coverage counts each premise once, which makes it a set rather than a rate and
is why Section~\ref{sec:coverage} reports it with no per-proof standardization;
it also makes coverage blind to usage. Counting the same premises once per
occurrence gives a different picture. The machine corpus resolves
\machineOccurrences{} premise occurrences from its \machineVocab{} distinct
premises, against \humanOccurrences{} from \humanVocab{} on the human side, and
it concentrates on a stratum essentially absent from mathlib: proofs closing in
at most two steps with an arithmetic-flavoured tactic set are
\degenerateMachineSharePct{}\% of the machine corpus and
\degenerateHumanSharePct{}\% of mathlib, \degenerateHumanCount{} of
\humanProofs{} declarations. We report that as a result rather than as sample
attrition. It is the largest cell of the machine corpus, it has no human
counterpart, and it is why the common-support design sets roughly a quarter of
the machine corpus aside. Usage concentrates to match. \texttt{sq\_nonneg}
alone accounts for \sqNonnegOccurrencePct{}\% of machine premise occurrences
and the twenty-five most frequent for \topTwentyFiveOccurrencePct{}\%, both
descriptive and post-hoc, and both not regenerable from the released artifact
(Section~\ref{sec:limitations}, L6). A machine proof resolves
\machinePremisesPerProof{} premises on average against
\humanPremisesPerProof{} for a human declaration. Coverage and concentration
are independent and neither substitutes for the other. Heavy reuse of a few
lemmas is ordinary in competition mathematics, and coverage alone would leave
open that a shared vocabulary is used comparably.

\section{\texttt{dag\_layer}: the reading under which it is correct}
\label{app:daglayer}

Section~\ref{sec:dag-layer} records the column as a centrality proxy and the
error in reading it as depth as ours. The attribute may be behaving as its
producers intended: if edges run from a declaration to its dependencies and the
computation seeds at in-degree-zero nodes, a column that puts uncited leaves at
layer zero and \texttt{Eq.refl} at the top is exactly what the definition
yields. Nothing here is a defect report against the release.

A second hazard is worth naming for other consumers. \layerTruncatedPct{}\% of
nodes carry a partial value that never finished propagating and is
indistinguishable in the published data from a completed one, so a consumer who
filters or bins on this column is mixing two populations without being told.

\section{What the coverage quantities do and do not read}
\label{app:construct}

These quantities read no edges. The objection is fair and we state it first.
Computing $J$ consults only membership, so coverage is a measure on node sets. The
graph supplies the definition of a node (which names are declarations, and
which layer each belongs to) rather than the arithmetic. The section headings retain the phrase \emph{structural
coverage}, and the reader should take it as naming the pair of measures defined
here rather than as a claim about the coverage arithmetic. Coverage is the
node-set family; breadth is the edge-derived one.

Two reporting conventions follow from that. We print the intersection count
beside every Jaccard so each index is recomputable from its own row, and we
give the human-only quantity as both a count and a fraction, because the two
vocabularies differ in size by orders of magnitude and a ratio alone would hide
which side is large. Coverage and breadth also answer different questions and
can move independently: coverage asks how much of the library a population
reaches, breadth how spread out that reach is, and a corpus may touch few
premises scattered across many communities or many premises within one.

The choice is deliberate rather than a convenience. Any edge-reading statistic
over these subgraphs requires aggregation: a centrality, a path length, a depth.
Aggregation over this graph is what Li et al.'s Finding~3 shows to be
uninterpretable unless the two hub layers are separated first, and
Section~\ref{sec:dag-layer} reports a second instance of the same trap on a
different quantity. A membership count needs no aggregation, so it needs no such
split to be read, and it survives controls that reverse or destroy the
edge-reading quantities.


\section*{NeurIPS Paper Checklist}

\begin{enumerate}

\item {\bf Claims}
    \item[] Question: Do the main claims made in the abstract and introduction accurately reflect the paper's contributions and scope?
    \item[] Answer: \answerYes{}
    \item[] Justification: The abstract and Section~\ref{sec:intro} state the three contributions the paper delivers: the instrument characterization of Section~\ref{sec:instrument}, the coverage results of Section~\ref{sec:results}, and the withdrawn finding of Section~\ref{sec:dag-layer}. Coverage is quoted as a range across designs rather than as a single figure, matching what Section~\ref{sec:results} reports.
    \item[] Guidelines:
    \begin{itemize}
        \item The answer \answerNA{} means that the abstract and introduction do not include the claims made in the paper.
        \item The abstract and/or introduction should clearly state the claims made, including the contributions made in the paper and important assumptions and limitations. A \answerNo{} or \answerNA{} answer to this question will not be perceived well by the reviewers. 
        \item The claims made should match theoretical and experimental results, and reflect how much the results can be expected to generalize to other settings. 
        \item It is fine to include aspirational goals as motivation as long as it is clear that these goals are not attained by the paper. 
    \end{itemize}

\item {\bf Limitations}
    \item[] Question: Does the paper discuss the limitations of the work performed by the authors?
    \item[] Answer: \answerYes{}
    \item[] Justification: Section~\ref{sec:limitations} is a dedicated limitations section. L1 states the instrument-transfer assumption as an assumption rather than a caveat and gives the direction of its residual error, L4 states the corpus-genre confound that no control available to us removes, and L6 states the single-corpus scope together with which reported figures are not reproducible from the artifact.
    \item[] Guidelines:
    \begin{itemize}
        \item The answer \answerNA{} means that the paper has no limitation while the answer \answerNo{} means that the paper has limitations, but those are not discussed in the paper. 
        \item The authors are encouraged to create a separate ``Limitations'' section in their paper.
        \item The paper should point out any strong assumptions and how robust the results are to violations of these assumptions (e.g., independence assumptions, noiseless settings, model well-specification, asymptotic approximations only holding locally). The authors should reflect on how these assumptions might be violated in practice and what the implications would be.
        \item The authors should reflect on the scope of the claims made, e.g., if the approach was only tested on a few datasets or with a few runs. In general, empirical results often depend on implicit assumptions, which should be articulated.
        \item The authors should reflect on the factors that influence the performance of the approach. For example, a facial recognition algorithm may perform poorly when image resolution is low or images are taken in low lighting. Or a speech-to-text system might not be used reliably to provide closed captions for online lectures because it fails to handle technical jargon.
        \item The authors should discuss the computational efficiency of the proposed algorithms and how they scale with dataset size.
        \item If applicable, the authors should discuss possible limitations of their approach to address problems of privacy and fairness.
        \item While the authors might fear that complete honesty about limitations might be used by reviewers as grounds for rejection, a worse outcome might be that reviewers discover limitations that aren't acknowledged in the paper. The authors should use their best judgment and recognize that individual actions in favor of transparency play an important role in developing norms that preserve the integrity of the community. Reviewers will be specifically instructed to not penalize honesty concerning limitations.
    \end{itemize}

\item {\bf Theory assumptions and proofs}
    \item[] Question: For each theoretical result, does the paper provide the full set of assumptions and a complete (and correct) proof?
    \item[] Answer: \answerNA{}
    \item[] Justification: The paper proves no theorems. Section~\ref{sec:definitions} gives measurement definitions rather than theoretical results, and each is defined constructively over the published dependency graph.
    \item[] Guidelines:
    \begin{itemize}
        \item The answer \answerNA{} means that the paper does not include theoretical results. 
        \item All the theorems, formulas, and proofs in the paper should be numbered and cross-referenced.
        \item All assumptions should be clearly stated or referenced in the statement of any theorems.
        \item The proofs can either appear in the main paper or the supplemental material, but if they appear in the supplemental material, the authors are encouraged to provide a short proof sketch to provide intuition. 
        \item Inversely, any informal proof provided in the core of the paper should be complemented by formal proofs provided in appendix or supplemental material.
        \item Theorems and Lemmas that the proof relies upon should be properly referenced. 
    \end{itemize}

    \item {\bf Experimental result reproducibility}
    \item[] Question: Does the paper fully disclose all the information needed to reproduce the main experimental results of the paper to the extent that it affects the main claims and/or conclusions of the paper (regardless of whether the code and data are provided or not)?
    \item[] Answer: \answerYes{}
    \item[] Justification: Section~\ref{sec:repro} gives the pinned graph revision, the pinned mathlib commit, the pinned machine-corpus revision and the single fixed seed, and Sections~\ref{sec:extraction} to \ref{sec:length-axis} and Section~\ref{sec:definitions} specify the extractor and every metric. Reproducibility is not uniform and the paper says so: most figures come from committed result files, a few are the graph release's own published figures and are obtained by re-fetching it at its pin, and a few are carried from the prior version of this work and are not reproducible from the artifact, which Section~\ref{sec:limitations}, L6, names.
    \item[] Guidelines:
    \begin{itemize}
        \item The answer \answerNA{} means that the paper does not include experiments.
        \item If the paper includes experiments, a \answerNo{} answer to this question will not be perceived well by the reviewers: Making the paper reproducible is important, regardless of whether the code and data are provided or not.
        \item If the contribution is a dataset and\slash or model, the authors should describe the steps taken to make their results reproducible or verifiable. 
        \item Depending on the contribution, reproducibility can be accomplished in various ways. For example, if the contribution is a novel architecture, describing the architecture fully might suffice, or if the contribution is a specific model and empirical evaluation, it may be necessary to either make it possible for others to replicate the model with the same dataset, or provide access to the model. In general. releasing code and data is often one good way to accomplish this, but reproducibility can also be provided via detailed instructions for how to replicate the results, access to a hosted model (e.g., in the case of a large language model), releasing of a model checkpoint, or other means that are appropriate to the research performed.
        \item While NeurIPS does not require releasing code, the conference does require all submissions to provide some reasonable avenue for reproducibility, which may depend on the nature of the contribution. For example
        \begin{enumerate}
            \item If the contribution is primarily a new algorithm, the paper should make it clear how to reproduce that algorithm.
            \item If the contribution is primarily a new model architecture, the paper should describe the architecture clearly and fully.
            \item If the contribution is a new model (e.g., a large language model), then there should either be a way to access this model for reproducing the results or a way to reproduce the model (e.g., with an open-source dataset or instructions for how to construct the dataset).
            \item We recognize that reproducibility may be tricky in some cases, in which case authors are welcome to describe the particular way they provide for reproducibility. In the case of closed-source models, it may be that access to the model is limited in some way (e.g., to registered users), but it should be possible for other researchers to have some path to reproducing or verifying the results.
        \end{enumerate}
    \end{itemize}


\item {\bf Open access to data and code}
    \item[] Question: Does the paper provide open access to the data and code, with sufficient instructions to faithfully reproduce the main experimental results, as described in supplemental material?
    \item[] Answer: \answerYes{}
    \item[] Justification: Section~\ref{sec:repro} describes the anonymized supplement: the extraction and analysis source, the committed result files behind the tables, the fetch scripts that obtain the input data at pinned revisions, and the test suite. Input data is not redistributed but is fetched from its public releases, and the figures no script in the artifact regenerates are named in Section~\ref{sec:limitations}, L6.
    \item[] Guidelines:
    \begin{itemize}
        \item The answer \answerNA{} means that paper does not include experiments requiring code.
        \item Please see the NeurIPS code and data submission guidelines (\url{https://neurips.cc/public/guides/CodeSubmissionPolicy}) for more details.
        \item While we encourage the release of code and data, we understand that this might not be possible, so \answerNo{} is an acceptable answer. Papers cannot be rejected simply for not including code, unless this is central to the contribution (e.g., for a new open-source benchmark).
        \item The instructions should contain the exact command and environment needed to run to reproduce the results. See the NeurIPS code and data submission guidelines (\url{https://neurips.cc/public/guides/CodeSubmissionPolicy}) for more details.
        \item The authors should provide instructions on data access and preparation, including how to access the raw data, preprocessed data, intermediate data, and generated data, etc.
        \item The authors should provide scripts to reproduce all experimental results for the new proposed method and baselines. If only a subset of experiments are reproducible, they should state which ones are omitted from the script and why.
        \item At submission time, to preserve anonymity, the authors should release anonymized versions (if applicable).
        \item Providing as much information as possible in supplemental material (appended to the paper) is recommended, but including URLs to data and code is permitted.
    \end{itemize}


\item {\bf Experimental setting/details}
    \item[] Question: Does the paper specify all the training and test details (e.g., data splits, hyperparameters, how they were chosen, type of optimizer) necessary to understand the results?
    \item[] Answer: \answerYes{}
    \item[] Justification: No model is trained or invoked here, so there are no hyperparameters or optimizer settings. The analogous choices are stated: the extractor's scoping and resolution rules in Section~\ref{sec:extraction}, the metric definitions in Section~\ref{sec:definitions}, the premise exclusion policy in Section~\ref{sec:coverage}, and the construction of each control (rarefaction, common support, area restriction, and 1:1 class-and-step matching) in Sections~\ref{sec:coverage} to \ref{sec:size-controls}, with pins and seed in Section~\ref{sec:repro}.
    \item[] Guidelines:
    \begin{itemize}
        \item The answer \answerNA{} means that the paper does not include experiments.
        \item The experimental setting should be presented in the core of the paper to a level of detail that is necessary to appreciate the results and make sense of them.
        \item The full details can be provided either with the code, in appendix, or as supplemental material.
    \end{itemize}

\item {\bf Experiment statistical significance}
    \item[] Question: Does the paper report error bars suitably and correctly defined or other appropriate information about the statistical significance of the experiments?
    \item[] Answer: \answerYes{}
    \item[] Justification: Where a quantity carries sampling variation we report it: the instrument reference line in Section~\ref{sec:reference-line} is quoted with a bootstrap interval and with a band over the precedence orderings of the class partition, and Section~\ref{sec:size-controls} reports the spread across seeded rarefaction draws. The headline vocabulary counts are censuses of both released corpora and carry no sampling error; Section~\ref{sec:size-controls} also states where a single seeded draw is adequate and where it is not.
    \item[] Guidelines:
    \begin{itemize}
        \item The answer \answerNA{} means that the paper does not include experiments.
        \item The authors should answer \answerYes{} if the results are accompanied by error bars, confidence intervals, or statistical significance tests, at least for the experiments that support the main claims of the paper.
        \item The factors of variability that the error bars are capturing should be clearly stated (for example, train/test split, initialization, random drawing of some parameter, or overall run with given experimental conditions).
        \item The method for calculating the error bars should be explained (closed form formula, call to a library function, bootstrap, etc.)
        \item The assumptions made should be given (e.g., Normally distributed errors).
        \item It should be clear whether the error bar is the standard deviation or the standard error of the mean.
        \item It is OK to report 1-sigma error bars, but one should state it. The authors should preferably report a 2-sigma error bar than state that they have a 96\% CI, if the hypothesis of Normality of errors is not verified.
        \item For asymmetric distributions, the authors should be careful not to show in tables or figures symmetric error bars that would yield results that are out of range (e.g., negative error rates).
        \item If error bars are reported in tables or plots, the authors should explain in the text how they were calculated and reference the corresponding figures or tables in the text.
    \end{itemize}

\item {\bf Experiments compute resources}
    \item[] Question: For each experiment, does the paper provide sufficient information on the computer resources (type of compute workers, memory, time of execution) needed to reproduce the experiments?
    \item[] Answer: \answerNo{}
    \item[] Justification: The analysis is CPU-only and single-machine, with no model training or inference and no Lean invocation of our own (Section~\ref{sec:repro}), but the paper reports no machine type, memory figure or wall-clock time for any step.
    \item[] Guidelines:
    \begin{itemize}
        \item The answer \answerNA{} means that the paper does not include experiments.
        \item The paper should indicate the type of compute workers CPU or GPU, internal cluster, or cloud provider, including relevant memory and storage.
        \item The paper should provide the amount of compute required for each of the individual experimental runs as well as estimate the total compute. 
        \item The paper should disclose whether the full research project required more compute than the experiments reported in the paper (e.g., preliminary or failed experiments that didn't make it into the paper). 
    \end{itemize}
    
\item {\bf Code of ethics}
    \item[] Question: Does the research conducted in the paper conform, in every respect, with the NeurIPS Code of Ethics \url{https://neurips.cc/public/EthicsGuidelines}?
    \item[] Answer: \answerYes{}
    \item[] Justification: The work analyses already-public, machine-checked proof corpora and a published dependency graph. It involves no human subjects and no personal data, anonymity is preserved in both the paper and the supplement, and LLM assistance is disclosed in Section~\ref{sec:repro}.
    \item[] Guidelines:
    \begin{itemize}
        \item The answer \answerNA{} means that the authors have not reviewed the NeurIPS Code of Ethics.
        \item If the authors answer \answerNo, they should explain the special circumstances that require a deviation from the Code of Ethics.
        \item The authors should make sure to preserve anonymity (e.g., if there is a special consideration due to laws or regulations in their jurisdiction).
    \end{itemize}


\item {\bf Broader impacts}
    \item[] Question: Does the paper discuss both potential positive societal impacts and negative societal impacts of the work performed?
    \item[] Answer: \answerNA{}
    \item[] Justification: The paper is a measurement study over released, formally verified proof corpora. It introduces no model, no deployment, and no data about people. Its foreseeable effect is on how proof automation is evaluated, which Section~\ref{sec:intro} states directly, and we identify no path from the work to a societal harm.
    \item[] Guidelines:
    \begin{itemize}
        \item The answer \answerNA{} means that there is no societal impact of the work performed.
        \item If the authors answer \answerNA{} or \answerNo, they should explain why their work has no societal impact or why the paper does not address societal impact.
        \item Examples of negative societal impacts include potential malicious or unintended uses (e.g., disinformation, generating fake profiles, surveillance), fairness considerations (e.g., deployment of technologies that could make decisions that unfairly impact specific groups), privacy considerations, and security considerations.
        \item The conference expects that many papers will be foundational research and not tied to particular applications, let alone deployments. However, if there is a direct path to any negative applications, the authors should point it out. For example, it is legitimate to point out that an improvement in the quality of generative models could be used to generate Deepfakes for disinformation. On the other hand, it is not needed to point out that a generic algorithm for optimizing neural networks could enable people to train models that generate Deepfakes faster.
        \item The authors should consider possible harms that could arise when the technology is being used as intended and functioning correctly, harms that could arise when the technology is being used as intended but gives incorrect results, and harms following from (intentional or unintentional) misuse of the technology.
        \item If there are negative societal impacts, the authors could also discuss possible mitigation strategies (e.g., gated release of models, providing defenses in addition to attacks, mechanisms for monitoring misuse, mechanisms to monitor how a system learns from feedback over time, improving the efficiency and accessibility of ML).
    \end{itemize}
    
\item {\bf Safeguards}
    \item[] Question: Does the paper describe safeguards that have been put in place for responsible release of data or models that have a high risk for misuse (e.g., pre-trained language models, image generators, or scraped datasets)?
    \item[] Answer: \answerNA{}
    \item[] Justification: The artifact contains analysis code and derived result files computed over public Lean corpora. No model, no scraped data and no generative capability is released, so there is nothing carrying a misuse risk of the kind this item covers.
    \item[] Guidelines:
    \begin{itemize}
        \item The answer \answerNA{} means that the paper poses no such risks.
        \item Released models that have a high risk for misuse or dual-use should be released with necessary safeguards to allow for controlled use of the model, for example by requiring that users adhere to usage guidelines or restrictions to access the model or implementing safety filters. 
        \item Datasets that have been scraped from the Internet could pose safety risks. The authors should describe how they avoided releasing unsafe images.
        \item We recognize that providing effective safeguards is challenging, and many papers do not require this, but we encourage authors to take this into account and make a best faith effort.
    \end{itemize}

\item {\bf Licenses for existing assets}
    \item[] Question: Are the creators or original owners of assets (e.g., code, data, models), used in the paper, properly credited and are the license and terms of use explicitly mentioned and properly respected?
    \item[] Answer: \answerNo{}
    \item[] Justification: Every asset is credited by citation and pinned: the dependency graph release in Section~\ref{sec:background}, the machine proof corpus and the human reference corpus in Section~\ref{sec:related}, and the exact revisions in Section~\ref{sec:repro}. None of them is redistributed, since the artifact fetches rather than ships them. The paper does not, however, state the license under which each asset is published.
    \item[] Guidelines:
    \begin{itemize}
        \item The answer \answerNA{} means that the paper does not use existing assets.
        \item The authors should cite the original paper that produced the code package or dataset.
        \item The authors should state which version of the asset is used and, if possible, include a URL.
        \item The name of the license (e.g., CC-BY 4.0) should be included for each asset.
        \item For scraped data from a particular source (e.g., website), the copyright and terms of service of that source should be provided.
        \item If assets are released, the license, copyright information, and terms of use in the package should be provided. For popular datasets, \url{paperswithcode.com/datasets} has curated licenses for some datasets. Their licensing guide can help determine the license of a dataset.
        \item For existing datasets that are re-packaged, both the original license and the license of the derived asset (if it has changed) should be provided.
        \item If this information is not available online, the authors are encouraged to reach out to the asset's creators.
    \end{itemize}

\item {\bf New assets}
    \item[] Question: Are new assets introduced in the paper well documented and is the documentation provided alongside the assets?
    \item[] Answer: \answerYes{}
    \item[] Justification: The only new asset is the analysis artifact, and Section~\ref{sec:repro} documents what it contains, how its input data is obtained at pinned revisions, how the paper builds from it, and which reported figures it does not regenerate.
    \item[] Guidelines:
    \begin{itemize}
        \item The answer \answerNA{} means that the paper does not release new assets.
        \item Researchers should communicate the details of the dataset\slash code\slash model as part of their submissions via structured templates. This includes details about training, license, limitations, etc. 
        \item The paper should discuss whether and how consent was obtained from people whose asset is used.
        \item At submission time, remember to anonymize your assets (if applicable). You can either create an anonymized URL or include an anonymized zip file.
    \end{itemize}

\item {\bf Crowdsourcing and research with human subjects}
    \item[] Question: For crowdsourcing experiments and research with human subjects, does the paper include the full text of instructions given to participants and screenshots, if applicable, as well as details about compensation (if any)? 
    \item[] Answer: \answerNA{}
    \item[] Justification: The paper involves no crowdsourcing and no research with human subjects. Both corpora are pre-existing released artifacts.
    \item[] Guidelines:
    \begin{itemize}
        \item The answer \answerNA{} means that the paper does not involve crowdsourcing nor research with human subjects.
        \item Including this information in the supplemental material is fine, but if the main contribution of the paper involves human subjects, then as much detail as possible should be included in the main paper. 
        \item According to the NeurIPS Code of Ethics, workers involved in data collection, curation, or other labor should be paid at least the minimum wage in the country of the data collector. 
    \end{itemize}

\item {\bf Institutional review board (IRB) approvals or equivalent for research with human subjects}
    \item[] Question: Does the paper describe potential risks incurred by study participants, whether such risks were disclosed to the subjects, and whether Institutional Review Board (IRB) approvals (or an equivalent approval/review based on the requirements of your country or institution) were obtained?
    \item[] Answer: \answerNA{}
    \item[] Justification: The paper involves no research with human subjects, so no IRB approval or equivalent review was required.
    \item[] Guidelines:
    \begin{itemize}
        \item The answer \answerNA{} means that the paper does not involve crowdsourcing nor research with human subjects.
        \item Depending on the country in which research is conducted, IRB approval (or equivalent) may be required for any human subjects research. If you obtained IRB approval, you should clearly state this in the paper. 
        \item We recognize that the procedures for this may vary significantly between institutions and locations, and we expect authors to adhere to the NeurIPS Code of Ethics and the guidelines for their institution. 
        \item For initial submissions, do not include any information that would break anonymity (if applicable), such as the institution conducting the review.
    \end{itemize}

\item {\bf Declaration of LLM usage}
    \item[] Question: Does the paper describe the usage of LLMs if it is an important, original, or non-standard component of the core methods in this research? Note that if the LLM is used only for writing, editing, or formatting purposes and does \emph{not} impact the core methodology, scientific rigor, or originality of the research, declaration is not required.
    %this research? 
    \item[] Answer: \answerNA{}
    \item[] Justification: No LLM is a component of the method. Every measurement here is a deterministic computation over a published dependency graph and released proof corpora, and no model is invoked anywhere in the pipeline; the machine proof corpus is the object of study rather than part of the method, and was produced and verified by its publishers. LLM assistance in implementation, data processing, running experiments and drafting is separately disclosed in Section~\ref{sec:repro}, as this venue's call requires, a broader requirement than this item's.
    \item[] Guidelines:
    \begin{itemize}
        \item The answer \answerNA{} means that the core method development in this research does not involve LLMs as any important, original, or non-standard components.
        \item Please refer to our LLM policy in the NeurIPS handbook for what should or should not be described.
    \end{itemize}

\end{enumerate}

\end{document}
